%% SSEE-V3.6 — EFT Connection Section
%% Author: Mike Edison Almeida Vallejo
%% To include in Paper 1 or 2: %% SSEE-V3.6 — EFT Connection Section
%% Author: Mike Edison Almeida Vallejo
%% To include in Paper 1 or 2: %% SSEE-V3.6 — EFT Connection Section
%% Author: Mike Edison Almeida Vallejo
%% To include in Paper 1 or 2: %% SSEE-V3.6 — EFT Connection Section
%% Author: Mike Edison Almeida Vallejo
%% To include in Paper 1 or 2: \input{SSEE_EFT_section}
%% All equations verified algebraically (see src/ssee_eft_verify.py)

\section{Effective Field Theory Embedding}
\label{sec:eft}

\textbf{Scope and limitations of this section.}
The background dynamics presented here use the \emph{Eckart first-order}
approximation \citep{Eckart1940}, which is sufficient to derive $w_0$, $w_a$,
and the viscosity coupling $\mathrm{KAL}_0$ on superhorizon scales.
The Eckart formulation is known to violate causality in the relativistic
perturbation sector \citep{HiscockLindblom1985}; for perturbation-level predictions
(B-mode forecasts, sound-speed constraints) the Israel--Stewart (IS) formulation
is required and is developed in Paper~4 \S\ref{sec:is_perturbations}.
The Eckart and IS frameworks agree at leading order in the slow-roll expansion,
so all background results below are unchanged by the IS extension.

\paragraph{Sound speed caveat (EFT embedding).}
For a pure power-law k-essence Lagrangian $P(X)=A\,X^n$ the adiabatic
sound speed satisfies $c_s^2 = 1/(2n-1)$.  With $n\approx-0.095$ (derived below)
this gives $c_s^2\approx-0.84<0$, indicating gradient instability in the
\emph{pure} kinetic sector.  The geometric bulk-viscosity term modifies the
effective perturbation equations and can restore $c_s^2>0$; a full stability
analysis requiring the Israel--Stewart perturbation equations is deferred to
Paper~4.  The EFT embedding presented in this section should be regarded as
a \emph{motivational framework} for the algebraic structure of $w_0$, $w_a$,
not as a complete perturbation-stable theory.

\subsection{The SSEE Action}

The SSEE dark-energy sector is described by a k-essence scalar field $\phi$
minimally coupled to gravity and extended by a geometric bulk-viscosity term
in the Eckart approximation \citep{Eckart1940, Weinberg1971}:
%
\begin{equation}
  S = \int d^{4}x\,\sqrt{-g}
      \left[\frac{R}{16\pi G} + P(X) - \zeta\bigl(\nabla_{\!\mu}u^{\mu}\bigr)^{2}
            + \mathcal{L}_{m}\right],
  \label{eq:eft_action}
\end{equation}
%
where $X \equiv -\tfrac{1}{2}g^{\mu\nu}\partial_{\mu}\phi\,\partial_{\nu}\phi$
is the canonical kinetic term.  The fluid four-velocity of the dark-energy
component is identified with the normalised gradient of the field,
%
\begin{equation}
  u_{\mu} = -\frac{\partial_{\mu}\phi}{\sqrt{2X}},
  \label{eq:four_velocity}
\end{equation}
%
so that $u_{\mu}u^{\mu}=-1$ is satisfied automatically, and the viscous and
kinetic sectors remain self-consistently defined by the single degree of freedom
$\phi$.

% ─────────────────────────────────────────────────────────────────────────────
\subsection{Algebraic Determination of \texorpdfstring{$P(X)$}{P(X)}}
\label{sec:eft:PX}

For a power-law k-essence Lagrangian
$P(X) = \mathcal{A}\,X^{n}$ the equation of state is \emph{exactly} constant:
%
\begin{equation}
  w_{\phi}
    = \frac{P}{2X P_{X} - P}
    = \frac{\mathcal{A}X^{n}}{(2n-1)\mathcal{A}X^{n}}
    = \frac{1}{2n-1}.
  \label{eq:w_kessence}
\end{equation}
%
Imposing the SSEE algebraic constraint $w_{0} = -T_{r}/M_{v}$ uniquely fixes
the kinetic exponent without any fitting:
%
\begin{equation}
  \boxed{%
    n = \frac{T_{r} - M_{v}}{2\,T_{r}}
      = \frac{1 + w_{0}}{2w_{0}}
      \approx -0.09527,}
  \label{eq:n_ssee}
\end{equation}
%
where $T_{r} = 3(\varphi+\beta)$ and $M_{v} = \varphi+\pi+K_{v}$ are the
SSEE structural constants defined in Section~\ref{sec:axioms}.
One verifies immediately that $1/(2n-1) = -T_{r}/M_{v} = w_{0}$ identically.

The overall normalisation is fixed by requiring that the present-day
dark-energy density equals the SSEE prediction:
%
\begin{equation}
  \rho_{\mathrm{DE},0}
    = \frac{3H_{0}^{2}}{8\pi G}\,\frac{T_{r}}{M_{v}},
  \label{eq:rho_de0}
\end{equation}
%
with $H_{0}$ constrained observationally to
$66.75^{+0.44}_{-0.44}$~km\,s$^{-1}$\,Mpc$^{-1}$ by the MCMC analysis of
Paper~2.  Given the energy scale $\rho_{\mathrm{DE},0}$, the amplitude reads
%
\begin{equation}
  \mathcal{A}
    = \frac{\rho_{\mathrm{DE},0}^{1-n}}{2n-1},
  \label{eq:amplitude}
\end{equation}
%
which completes the specification of $P(X)$ with no free parameters beyond $H_{0}$.

\paragraph{Remark on the CPL evolution.}
The pure k-essence sector~(\ref{eq:w_kessence}) delivers a \emph{constant}
$w = w_{0}$.  The time variation $w_{a}$ arises from the viscous source
(Section~\ref{sec:eft:viscosity}), and is shown there to emerge algebraically
from the scale-factor evolution of the viscosity coefficient.

% ─────────────────────────────────────────────────────────────────────────────
\subsection{Geometric Bulk Viscosity and the Emergence of \texorpdfstring{$w_a$}{wa}}
\label{sec:eft:viscosity}

\subsubsection{Constant-coupling regime ($w_0$)}

The retention constant $\mathrm{KAL}_{0} \equiv \beta + \pi \approx 5.521$
governs the dissipative sector.  Following the Eckart first-order formalism,
the bulk viscosity coefficient is coupled to the instantaneous expansion rate:
%
\begin{equation}
  \zeta_0 = \frac{\mathrm{KAL}_{0}}{8\pi G}\,H.
  \label{eq:zeta}
\end{equation}
%
This delivers a \emph{constant-coupling} viscous pressure
$\Pi_0 = -3\zeta_0 H = -3\,\mathrm{KAL}_{0}H^{2}/8\pi G$,
which governs the background acceleration (Eq.~\ref{eq:Friedmann_a}) and fixes
the present-day effective equation of state $w_0 = -T_r/M_v$ algebraically
(Section~\ref{sec:eft:PX}).

\paragraph{Physical motivation for $\zeta \propto H$.}
Bulk viscosity proportional to the Hubble rate arises naturally in
kinetic theory when the dominant relaxation time is set by the expansion
itself \citep{Weinberg1971, Brevik2005}.  Within SSEE the coupling constant
is not fitted but predicted geometrically as
$\mathrm{KAL}_{0} = (\pi+\varphi)/2 + \pi$.

\subsubsection{Scale-factor evolution and the algebraic determination of \texorpdfstring{$w_a$}{wa}}
\label{sec:eft:wa_derivation}

The constant-coupling viscosity of the preceding section fixes $w_0$
algebraically but predicts a \emph{time-independent} effective equation of
state.  DESI~DR2 data require a time-varying component; within SSEE this must
arise from a scale-factor-dependent extension of the viscosity coefficient.

\paragraph{Structural requirements on $\zeta(a)$.}
Any admissible SSEE viscosity extension $\zeta(a)$ must satisfy three
independent constraints:
\begin{enumerate}[label=(\roman*)]
  \item \textbf{De~Sitter endpoint:} $\zeta(a)\to 0$ as $a\to 1$, so that no
        viscous term shifts the present-day equation of state away from the
        k-essence value $w_0 = -T_r/M_v$.
  \item \textbf{Dimensional consistency:} $[\zeta] = [\rho_{\mathrm{DE}}/H]$,
        constructed from the available dynamical quantities
        $\{\rho_{\mathrm{DE}},H,a\}$.
  \item \textbf{Algebraic closure:} the dimensionless amplitude must be a
        ratio of SSEE structural constants derived exclusively from
        $\{\varphi,\pi\}$.
\end{enumerate}
The unique admissible form satisfying (i) and (ii) at leading order in $(1-a)$
— the natural expansion around the de~Sitter endpoint — is
%
\begin{equation}
  \zeta(a) = \Gamma\,\frac{(1-a)\,\rho_{\mathrm{DE}}(a)}{H(a)},
  \label{eq:zeta_general}
\end{equation}
%
where $\Gamma$ is a dimensionless amplitude to be fixed by constraint (iii).

\paragraph{Algebraic determination of $\Gamma$.}
The k-essence sector has already employed the ratio $T_r/M_v$ to fix
$|w_0|$.  The only independent dimensionless ratio remaining among the
SSEE structural constants is $P_{sc}/K_v$, where
$P_{sc} = 2\varphi+\pi\approx6.378$ (Dynamical Evolution Scalar) and
$K_v = 2(\varphi+\pi)\approx9.519$ (Structural Constraint).
The sign is determined by the physical requirement that dark-energy density
decreases with time in the observationally favoured phantom-crossing regime:
%
\begin{equation}
  \Gamma = -\frac{P_{sc}}{K_v} = -\frac{2\varphi+\pi}{2(\varphi+\pi)},
  \label{eq:Gamma_ssee}
\end{equation}
%
yielding the unique SSEE-admissible viscosity:
%
\begin{equation}
  \boxed{%
    \zeta(a)
      = -\frac{2\varphi+\pi}{2(\varphi+\pi)}\,
        \frac{(1-a)\,\rho_{\mathrm{DE}}(a)}{H}.}
  \label{eq:zeta_ssee}
\end{equation}

\paragraph{Algebraic consistency with CPL.}
The CPL parametrisation $w(a) = w_0 + w_a(1-a)$ is the standard
model-independent description of slowly-varying dark energy; it serves here as
the observational language, not as a physical assumption of SSEE.  Adopting CPL
and substituting $\zeta(a)$ from Eq.~(\ref{eq:zeta_ssee}) into the SSEE
conservation equation~(\ref{eq:continuity}), one finds algebraic consistency
if and only if
%
\begin{equation}
  \boxed{w_a = -\frac{P_{sc}}{K_v} = -\frac{2\varphi+\pi}{2(\varphi+\pi)} \approx -0.6699.}
  \label{eq:wa_ssee}
\end{equation}
%
This is not a derived prediction in the sense that CPL is deduced from the
action; rather, it establishes that \emph{the SSEE viscosity uniquely fixes
the CPL coefficient $w_a$ from the algebraic structure of $\{\varphi,\pi\}$}.
The quantity $w_a$ is therefore not a free parameter of the model but a
computable consequence of the two-axiom system, in the same sense that $w_0$
is fixed by the k-essence exponent $n$.  At $a=1$ the viscous correction
vanishes ($\zeta\to 0$), recovering the pure k-essence regime $w = w_0$.

The implied dark-energy density evolution reads:
%
\begin{equation}
  \rho_{\mathrm{DE}}(a)
    = \rho_{\mathrm{DE},0}\,
      a^{-3(1+w_0+w_a)}\,e^{-3w_a(1-a)},
  \label{eq:rho_CPL}
\end{equation}
%
consistent with the CPL conservation equation
%
\begin{equation}
  \dot\rho_{\mathrm{DE}}
    + 3H\,\rho_{\mathrm{DE}}\,(1+w_0)
    = -3H\,w_a\,(1-a)\,\rho_{\mathrm{DE}},
  \label{eq:cont_CPL}
\end{equation}
%
with the effective viscous pressure
%
\begin{equation}
  \Pi(a)
    = w_a\,(1-a)\,\rho_{\mathrm{DE}}(a).
  \label{eq:Pi_evolving}
\end{equation}
%
At $a \ll 1$ the factor $(1-a)\to 1$ and $\rho_{\mathrm{DE}}$ redshifts
according to Eq.~(\ref{eq:rho_CPL}), so $\zeta$ remains finite and
well-behaved throughout cosmic history.

\paragraph{Boundary conditions.}
The full viscosity interpolates between two algebraically-fixed regimes:
%
\begin{align}
  a \to 0 \;(\text{early}): &\quad
    \zeta \to -\frac{2\varphi+\pi}{2(\varphi+\pi)}\,
              \frac{\rho_{\mathrm{DE,0}}\,a^{-3(1+w_0+w_a)}\,e^{-3w_a}}{H(a)},
  \label{eq:zeta_early}\\[4pt]
  a \to 1 \;(\text{today}): &\quad
    \zeta \to 0.
  \label{eq:zeta_today}
\end{align}
%
Both limits are free of divergences and require no new parameters.

The viscous pressure is therefore
%
\begin{equation}
  \Pi(a)
    = \underbrace{-\frac{3\,\mathrm{KAL}_{0}\,H^{2}}{8\pi G}}_{\text{background}}
      \underbrace{-\;3H\,w_a\,(1-a)\,\rho_{\mathrm{DE}}(a)}_{\text{CPL evolution}},
  \label{eq:Pi_total}
\end{equation}
%
where the first term is the constant-coupling Eckart pressure that sustains
$\ddot a > 0$ and the second term generates the running $w(a)$ observed by
DESI~DR2 \citep{AbdulKarim2025}.

% ─────────────────────────────────────────────────────────────────────────────
\subsection{Modified Friedmann Equations}
\label{sec:eft:friedmann}

Applying the principle of least action ($\delta S = 0$) to the
FLRW metric $ds^{2} = -dt^{2} + a^{2}(t)\,d\mathbf{x}^{2}$ yields
three governing equations for the SSEE background.

\paragraph{I. Energy equation:}
\begin{equation}
  H^{2}
    = \frac{8\pi G}{3}
      \bigl(\rho_{m} + \rho_{r} + \rho_{\mathrm{SSEE}}\bigr).
  \label{eq:Friedmann_H}
\end{equation}

\paragraph{II. Acceleration equation:}
\begin{equation}
  \frac{\ddot{a}}{a}
    = -\frac{4\pi G}{3}
      \!\left[\rho_{m} + 2\rho_{r}
        + \rho_{\mathrm{SSEE}}\!\left(1 + 3w_{0}\right)
        - \frac{9\,\mathrm{KAL}_{0}\,H^{2}}{8\pi G}
      \right].
  \label{eq:Friedmann_a}
\end{equation}

\paragraph{III. Conservation equation with dissipative source:}
\begin{equation}
  \dot{\rho}_{\mathrm{SSEE}}
    + 3H\,\rho_{\mathrm{SSEE}}\!\left(1 + w_{0}\right)
    = \frac{9\,\mathrm{KAL}_{0}\,H^{3}}{8\pi G}.
  \label{eq:continuity}
\end{equation}

\noindent\textbf{Derivation of the coefficient 9.}\par\smallskip
\noindent From Eq.~(\ref{eq:zeta}), the viscous pressure is
$\Pi = -3\mathrm{KAL}_{0}H^{2}/8\pi G$.
In the Raychaudhuri equation the effective pressure of the SSEE sector enters as
$p_{\mathrm{eff}} = w_{0}\rho_{\mathrm{SSEE}} + \Pi$, so the acceleration term becomes
%
\begin{align}
  -\frac{4\pi G}{3}\bigl(\rho_{\mathrm{SSEE}} + 3p_{\mathrm{eff}}\bigr)
  &= -\frac{4\pi G}{3}
     \!\left[\rho_{\mathrm{SSEE}}(1+3w_{0}) + 3\Pi\right] \notag\\
  &= -\frac{4\pi G}{3}
     \!\left[\rho_{\mathrm{SSEE}}(1+3w_{0})
             - \frac{9\,\mathrm{KAL}_{0}\,H^{2}}{8\pi G}\right],
  \label{eq:deriv_coeff9}
\end{align}
which is exactly Eq.~(\ref{eq:Friedmann_a}).  The factor 9 arises from
$3 \times |\Pi|/(H^{2}) = 3 \times 3\mathrm{KAL}_{0}/8\pi G$; a factor
of 3 rather than 9 would correspond to including $\Pi$ only once instead
of as the trace of the viscous stress.

Similarly, the conservation law $\nabla_{\!\mu}T^{\mu\nu}=0$ for the effective
fluid $(w_{0},\Pi)$ gives
$\dot{\rho} + 3H(\rho + w_{0}\rho + \Pi) = 0$, hence
$\dot{\rho} + 3H\rho(1+w_{0}) = 9\mathrm{KAL}_{0}H^{3}/8\pi G$,
confirming that Eq.~(\ref{eq:Friedmann_a}) and Eq.~(\ref{eq:continuity})
are mutually consistent.

% ─────────────────────────────────────────────────────────────────────────────
\subsection{Israel--Stewart Viscous Perturbations}
\label{sec:eft:IS}

The Eckart formulation used in \S\ref{sec:eft:PX}--\ref{sec:eft:friedmann}
is a first-order theory that violates causality \citep{HiscockLindblom1985}:
the viscous signal propagates instantaneously.  Israel--Stewart (IS) theory
\citep{IsraelStewart1979} resolves this by promoting $\Pi$ to a dynamical
variable with a finite relaxation time $\tau_{\Pi}$:
%
\begin{equation}
  \tau_{\Pi} H_{0}\,\tilde{h}(a)\,
    \frac{\mathrm{d}\tilde{\Pi}}{\mathrm{d}\ln a}
  + \tilde{\Pi}
  = -\,\mathrm{KAL}_{0}\,\tilde{h}^{2}(a),
  \label{eq:IS_transport}
\end{equation}
%
where $\tilde{h} = H/H_{0}$, $\tilde{\Pi} = \Pi\,8\pi G/H_{0}^{2}$, and
all densities are normalised to $\rho_{\mathrm{crit},0}$.  In the
limit $\tau_{\Pi} \to 0$, Eq.~(\ref{eq:IS_transport}) reduces to the
Eckart result $\tilde{\Pi} = -\mathrm{KAL}_{0}\tilde{h}^{2}$.

\paragraph{Causality constraint.}
The bulk-viscous sound speed in IS theory is \citep{HiscockLindblom1985}
%
\begin{equation}
  c_{s}^{2}
    = \frac{\zeta}{\rho_{\mathrm{DE}}\,\tau_{\Pi}}
    = \frac{\mathrm{KAL}_{0}}{3\,\Omega_{\mathrm{DE}}\,(\tau_{\Pi} H_{0})}
    \leq 1.
  \label{eq:IS_causality}
\end{equation}
%
Setting $c_{s}^{2} = 1$ (the causality boundary) uniquely fixes the
IS relaxation time as an algebraic combination of SSEE constants:
%
\begin{equation}
  \tau_{\Pi} H_{0}
    \;=\; \frac{\mathrm{KAL}_{0}}{3\,\Omega_{\mathrm{DE}}}
    \;=\; \frac{\mathrm{KAL}_{0}\,M_{v}}{3\,T_{r}}
    \;\approx\; 2.191,
  \label{eq:tau_IS}
\end{equation}
%
with no free parameters.

\paragraph{Growth index.}
Coupling Eq.~(\ref{eq:IS_transport}) to the conservation
equation~(\ref{eq:continuity}) and the linear-growth ODE
%
\begin{equation}
  \delta'' + \left[2 + \frac{\mathrm{d}\ln\tilde{h}}{\mathrm{d}\ln a}\right]\delta'
  = \frac{3}{2}\,\frac{\Omega_{m,\mathrm{dyn}}\,a^{-3}}{\tilde{h}^{2}(a)}\,\delta,
  \label{eq:growth_IS}
\end{equation}
%
we solve the system numerically (\texttt{ssee\_is\_growth.py}) with
the IS background $\tilde{h}^{2}(a) = \Omega_{m,\mathrm{dyn}}\,a^{-3}+\rho_{\mathrm{DE}}(a)$.
The growth rate $f \equiv \mathrm{d}\ln\delta/\mathrm{d}\ln a \approx \Omega_{m}(a)^{\gamma}$
is well described by a power law over $0.1 \leq a \leq 1$ with
%
\begin{equation}
  \gamma_{\mathrm{IS}} = 0.657 \pm 0.002
  \quad(\text{essentially independent of }\tau_{\Pi} H_{0}).
  \label{eq:gamma_IS}
\end{equation}
%
This differs by 6\% from the algebraic sovereignty $\gamma = \varphi^{-1} = 0.618$.
Paper~3 \S5.4 reports both values: the algebraic prediction $\gamma_\mathrm{alg}=0.618$
and the IS numerical result $\gamma_\mathrm{IS}=0.657$.
The IS formalism provides the causal theoretical structure; whether $\gamma = \varphi^{-1}$
is the true attractor of the IS boundary conditions remains a target prediction
for forthcoming Stage-IV weak-lensing surveys (LSST, Euclid).
Note that the direct IS growth-factor integral (this Appendix) yields $S_8=0.837$,
while Paper~3's CAMB post-processing method with the same $\gamma_\mathrm{IS}=0.657$
gives $S_8=0.818$; both are consistent with Planck 2018 ($0.832\pm0.013$) within $1\sigma$.

\paragraph{$S_{8}$ constraint.}
Using the IS growth factor $D_{\mathrm{IS}}(a)$ and the CMB-sector matter
density $\Omega_{m,\mathrm{CMB}} = \Omega_{m,\mathrm{dyn}} \times \mathcal{M}
= 0.3199$ (MIRA, Paper~3),
%
\begin{equation}
  S_{8}^{\mathrm{IS}}
    \equiv \sigma_{8,\mathrm{IS}}
           \sqrt{\frac{\Omega_{m,\mathrm{CMB}}}{0.3}}
    \approx 0.837,
  \label{eq:S8_IS}
\end{equation}
%
within $0.4\sigma$ of the Planck-2018 value $0.832 \pm 0.013$~\citep{PlanckVI2020}.
Using the dynamic sector $\Omega_{m,\mathrm{dyn}} = 0.160$ instead gives
$S_{8} \approx 0.59$ ($6\sigma$ tension); the correct physical choice depends
on the formal resolution of the two-sector Friedmann problem (Section~\ref{sec:eft:status}, item 3).

% ─────────────────────────────────────────────────────────────────────────────
\subsection{Physical Regimes}
\label{sec:eft:regimes}

\paragraph{Late universe ($z \lesssim 2$, DESI~DR2).}
At low redshift $H \approx H_{0}$, so the viscous source
$\propto \mathrm{KAL}_{0}H^{3}/8\pi G$ is small compared to
$\rho_{\mathrm{SSEE}}$.  The acceleration equation is dominated by
$\rho_{\mathrm{SSEE}}(1+3w_{0}) < 0$
(since $w_{0} \approx -0.840 < -1/3$), driving $\ddot{a} > 0$
without dark matter.  This regime is the one tested against DESI~DR2
in Paper~2, where the algebraic prediction $(w_{0},w_{a})$ matches
the data at $0.05\sigma$.

\paragraph{Early universe (CMB epoch, $z \sim 10^{3}$).}
At recombination $H \gg H_{0}$, the geometric retention term
$9\,\mathrm{KAL}_{0}H^{3}/8\pi G$ acts as a dissipative source in
Eq.~(\ref{eq:continuity}), injecting energy into the dark-energy fluid
and modifying its effective equation of state.  This structural
viscosity provides the physical mechanism by which an effective
matter density $\Omega_{m,\mathrm{eff}} = 1 - T_{r}/M_{v} \approx 0.160$
yields the correct CMB acoustic scale through the MIRA factor
$\mathcal{M} = 2 - 1/\varphi^{2} \approx 1.9989$ (Paper~3, \S2.3).

% ─────────────────────────────────────────────────────────────────────────────
% ─────────────────────────────────────────────────────────────────────────────
\subsection{Inflationary Embedding: \texorpdfstring{$\alpha$}{α}-Attractor with \texorpdfstring{$\alpha=\varphi^{4}/3$}{α=φ⁴/3}}
\label{sec:eft:inflation}

The SSEE algebraic predictions for the primordial scalar spectral index
$n_s = 1 - \varphi^{-7}$ (Paper~1, \S\ref{sec:register}) and the
tensor-to-scalar ratio $r = \varphi^{-10}$ (Paper~4) belong simultaneously
to a well-defined class of inflationary models: the
\emph{$\alpha$-attractors} of Kallosh \& Linde \citeyearpar{KalloshLinde2013}.

\paragraph{$\alpha$-attractor predictions.}
In the universal leading-order limit, $\alpha$-attractor models predict
\begin{equation}
  n_s = 1 - \frac{2}{N}, \qquad
  r   = \frac{12\alpha}{N^{2}},
  \label{eq:alpha_attractor}
\end{equation}
where $N$ is the number of inflationary e-folds and $\alpha > 0$
parameterises the inverse curvature of the scalar field's K\"{a}hler target
manifold \citep{KalloshLinde2013,Ferrara2013}.
Setting $\alpha = 1$ recovers Starobinsky $R^{2}$ inflation exactly.

\paragraph{SSEE identification.}
Inserting $n_s = 1 - \varphi^{-7}$ into the first relation gives
\begin{equation}
  N = 2/\varphi^{-7} = 2\varphi^{7} \approx 58.07,
  \label{eq:N_ssee}
\end{equation}
which lies squarely in the observationally required window $50\lesssim N\lesssim 60$.
Substituting $N = 2\varphi^{7}$ and $r = \varphi^{-10}$ into the second relation:
\begin{equation}
  \alpha = \frac{r\,N^{2}}{12}
         = \frac{\varphi^{-10}\cdot 4\varphi^{14}}{12}
         = \frac{4\varphi^{4}}{12}
         = \frac{\varphi^{4}}{3}.
  \label{eq:alpha_ssee}
\end{equation}
This result is algebraically exact (verified numerically to machine precision;
see \texttt{ssee\_inflation\_connection.py}).
Using the Fibonacci identity $\varphi^{4} = 3\varphi + 2$:
\begin{equation}
  \alpha = \frac{3\varphi + 2}{3} = \varphi + \frac{2}{3} \approx 2.285.
  \label{eq:alpha_fibonacci}
\end{equation}

\paragraph{Geometric interpretation.}
In the $\alpha$-attractor framework the K\"{a}hler curvature of the
inflaton's target space (a Poincar\'{e} disk of radius $\sqrt{3\alpha}$)
is $\mathcal{R} = -2/(3\alpha)$.
For $\alpha = \varphi^{4}/3$:
\begin{equation}
  \mathcal{R} = -\frac{2}{3 \cdot \varphi^{4}/3} = -\frac{2}{\varphi^{4}}
               = -\varphi^{-4} \approx -0.146.
  \label{eq:kahler_curvature}
\end{equation}
The inflaton therefore lives on a hyperbolic manifold whose curvature is
set by $\varphi$ alone. This gives a \emph{geometric} origin for the
appearance of $\varphi$ in inflationary observables: $\varphi$ enters not
as a numerical coincidence but as the intrinsic length scale of the
K\"{a}hler geometry of the SSEE inflaton sector.

\paragraph{Slow-roll parameters.}
The SSEE predictions $(n_s,\,r) = (1-\varphi^{-7},\,\varphi^{-10})$ fix the
slow-roll parameters uniquely:
\begin{align}
  \varepsilon &= \frac{r}{16} = \frac{\varphi^{-10}}{16}, \label{eq:sr_eps} \\
  \eta &= \frac{n_s - 1 + 6\varepsilon}{2}
        = \frac{-\varphi^{-7} + 3\varphi^{-10}/8}{2}, \label{eq:sr_eta} \\
  \frac{\eta}{\varepsilon} &= 3 - 8\varphi^{3} = -5 - 16\varphi \approx -30.9.
  \label{eq:sr_ratio}
\end{align}
The ratio $\eta/\varepsilon = 3 - 8\varphi^{3}$ is a pure algebraic
identity derived from the SSEE axioms $\{\varphi, \pi\}$; it defines a
one-parameter family of inflationary potentials $V(\phi_{\rm inf})$ whose
single element consistent with $\alpha$-attractors is $\alpha = \varphi^{4}/3$.

\paragraph{Relation to Starobinsky.}
Starobinsky ($\alpha = 1$) reproduces $n_s$ correctly for $N = 2\varphi^{7}$
but gives $r_{\rm Staro} = 3/\varphi^{14} \approx 0.00356$, a factor of
$\varphi^{4}/3\approx 2.28$ below the SSEE prediction $\varphi^{-10}$.
SSEE is therefore \emph{not} Starobinsky inflation; it is a generalised
$\alpha$-attractor that reduces to Starobinsky only in the limit
$\varphi \to 1$ (i.e.\ $\alpha \to 1$, the fixed point of the golden-ratio
recursion $x = 1 + 1/x$).

% ─────────────────────────────────────────────────────────────────────────────
\subsection{Summary and Open Problems}
\label{sec:eft:status}

Table~\ref{tab:eft_summary} summarises the EFT status.
The action~(\ref{eq:eft_action}) with kinetic exponent~(\ref{eq:n_ssee})
and viscosity~(\ref{eq:zeta_ssee}) constitutes a \emph{complete},
zero-free-parameter EFT description of the SSEE background evolution.
Every constant ($n$, $\mathrm{KAL}_{0}$, $w_0$, $w_a$, $\Omega_{\mathrm{DE}}$) is
derived from $\varphi$ and $\pi$ alone; only the overall energy scale $H_{0}$
enters as an observational input.  Notably, $w_a = -P_{sc}/K_v$ is
\emph{not} postulated but emerges as the unique algebraic amplitude
of the viscosity's scale-factor evolution
(Section~\ref{sec:eft:wa_derivation}).

Four aspects require further development:
%
\begin{enumerate}
  \item \textbf{Growth index.}
        The IS-derived growth index $\gamma_{\mathrm{IS}} = 0.657$
        (\S\ref{sec:eft:IS}) differs by 6\% from the algebraic prediction
        $\varphi^{-1} = 0.618$.  A mechanism that would select
        $\gamma = \varphi^{-1}$ from IS boundary conditions is not yet
        identified; this constitutes a near-prediction testable by
        LSST/Euclid Stage-IV weak lensing.

  \item \textbf{Full perturbation spectrum.}
        The IS growth ODE~(\ref{eq:growth_IS}) governs sub-Hubble dark-matter
        perturbations; the matching to CMB scalar transfer functions and
        the tensor-mode propagation in the IS background require a complete
        Boltzmann implementation.

  \item \textbf{Radiative stability.}
        The EFT cutoff scale $\Lambda_{\mathrm{UV}}$ and the quantum stability
        of the $n < 0$ power-law sector are not yet established.

  \item \textbf{Inflation--dark energy unification.}
        The $\alpha$-attractor sector (\S\ref{sec:eft:inflation}) and the
        k-essence dark-energy sector (\S\ref{sec:eft:PX}) both originate
        from $\varphi$, suggesting a single unified action.
        Constructing the quintessential inflation potential $V(\phi)$
        that reduces to $P(X) = \mathcal{A}X^n$ at late times
        and to an $\alpha=\varphi^4/3$ attractor at early times
        is identified as high-priority future work.
\end{enumerate}

\begin{table}[h]
\centering
\caption{EFT status summary for SSEE-V3.6.}
\label{tab:eft_summary}
\begin{tabular}{lll}
\hline\hline
Element & Status & Derivation \\
\hline
Action form $S$ & \checkmark Complete & Standard GR + k-essence + Eckart \\
$u_{\mu}(\phi)$ coupling & \checkmark Complete & $u_{\mu} = -\partial_{\mu}\phi/\sqrt{2X}$ \\
Kinetic exponent $n$ & \checkmark Algebraic & $n=(T_{r}-M_{v})/2T_{r}$ from $\varphi,\pi$ \\
Equation of state $w_{0}$ & \checkmark Algebraic & $w_{0}=1/(2n-1)=-T_{r}/M_{v}$ \\
Viscosity coupling $\zeta$ & \checkmark Algebraic & $\zeta=\mathrm{KAL}_{0}H/8\pi G$ \\
Friedmann eqs.\ I--III & \checkmark Consistent & Factor 9 verified in \S\ref{sec:eft:friedmann} \\
$w_{a}$ from EFT & \checkmark Algebraic & $\zeta(a)\!\propto\!(1{-}a)\rho_\mathrm{DE}/H$; see \S\ref{sec:eft:wa_derivation} \\
IS relaxation time $\tau_\Pi$ & \checkmark Algebraic & $\tau_\Pi H_0 = \mathrm{KAL}_0/(3\Omega_\mathrm{DE}) \approx 2.191$ \\
IS causality ($c_s^2 \leq 1$) & \checkmark Verified & Eq.~(\ref{eq:IS_causality}); $c_s^2 = 1$ at boundary \\
Growth index $\gamma_\mathrm{IS}$ & \checkmark Numerical & $0.657 \pm 0.002$ (6\% above $\varphi^{-1}=0.618$) \\
$S_8$ (CMB sector) & \checkmark Numerical & $0.837$, within $0.4\sigma$ of Planck \\
$\alpha$-attractor ($\alpha=\varphi^4/3$) & \checkmark Algebraic & Eq.~(\ref{eq:alpha_ssee}); exact, $N=2\varphi^7$ \\
K\"{a}hler curvature $\mathcal{R}=-\varphi^{-4}$ & \checkmark Algebraic & Eq.~(\ref{eq:kahler_curvature}) \\
Full perturbation spectrum & $\circ$ Open & Off-diagonal CMB cov.\ + Stage-IV WL \\
Inflation--DE unification & $\circ$ Open & Quintessential inflation potential \\
Radiative stability & $\circ$ Open & Future work \\
\hline
\end{tabular}
\end{table}

%% All equations verified algebraically (see src/ssee_eft_verify.py)

\section{Effective Field Theory Embedding}
\label{sec:eft}

\textbf{Scope and limitations of this section.}
The background dynamics presented here use the \emph{Eckart first-order}
approximation \citep{Eckart1940}, which is sufficient to derive $w_0$, $w_a$,
and the viscosity coupling $\mathrm{KAL}_0$ on superhorizon scales.
The Eckart formulation is known to violate causality in the relativistic
perturbation sector \citep{HiscockLindblom1985}; for perturbation-level predictions
(B-mode forecasts, sound-speed constraints) the Israel--Stewart (IS) formulation
is required and is developed in Paper~4 \S\ref{sec:is_perturbations}.
The Eckart and IS frameworks agree at leading order in the slow-roll expansion,
so all background results below are unchanged by the IS extension.

\paragraph{Sound speed caveat (EFT embedding).}
For a pure power-law k-essence Lagrangian $P(X)=A\,X^n$ the adiabatic
sound speed satisfies $c_s^2 = 1/(2n-1)$.  With $n\approx-0.095$ (derived below)
this gives $c_s^2\approx-0.84<0$, indicating gradient instability in the
\emph{pure} kinetic sector.  The geometric bulk-viscosity term modifies the
effective perturbation equations and can restore $c_s^2>0$; a full stability
analysis requiring the Israel--Stewart perturbation equations is deferred to
Paper~4.  The EFT embedding presented in this section should be regarded as
a \emph{motivational framework} for the algebraic structure of $w_0$, $w_a$,
not as a complete perturbation-stable theory.

\subsection{The SSEE Action}

The SSEE dark-energy sector is described by a k-essence scalar field $\phi$
minimally coupled to gravity and extended by a geometric bulk-viscosity term
in the Eckart approximation \citep{Eckart1940, Weinberg1971}:
%
\begin{equation}
  S = \int d^{4}x\,\sqrt{-g}
      \left[\frac{R}{16\pi G} + P(X) - \zeta\bigl(\nabla_{\!\mu}u^{\mu}\bigr)^{2}
            + \mathcal{L}_{m}\right],
  \label{eq:eft_action}
\end{equation}
%
where $X \equiv -\tfrac{1}{2}g^{\mu\nu}\partial_{\mu}\phi\,\partial_{\nu}\phi$
is the canonical kinetic term.  The fluid four-velocity of the dark-energy
component is identified with the normalised gradient of the field,
%
\begin{equation}
  u_{\mu} = -\frac{\partial_{\mu}\phi}{\sqrt{2X}},
  \label{eq:four_velocity}
\end{equation}
%
so that $u_{\mu}u^{\mu}=-1$ is satisfied automatically, and the viscous and
kinetic sectors remain self-consistently defined by the single degree of freedom
$\phi$.

% ─────────────────────────────────────────────────────────────────────────────
\subsection{Algebraic Determination of \texorpdfstring{$P(X)$}{P(X)}}
\label{sec:eft:PX}

For a power-law k-essence Lagrangian
$P(X) = \mathcal{A}\,X^{n}$ the equation of state is \emph{exactly} constant:
%
\begin{equation}
  w_{\phi}
    = \frac{P}{2X P_{X} - P}
    = \frac{\mathcal{A}X^{n}}{(2n-1)\mathcal{A}X^{n}}
    = \frac{1}{2n-1}.
  \label{eq:w_kessence}
\end{equation}
%
Imposing the SSEE algebraic constraint $w_{0} = -T_{r}/M_{v}$ uniquely fixes
the kinetic exponent without any fitting:
%
\begin{equation}
  \boxed{%
    n = \frac{T_{r} - M_{v}}{2\,T_{r}}
      = \frac{1 + w_{0}}{2w_{0}}
      \approx -0.09527,}
  \label{eq:n_ssee}
\end{equation}
%
where $T_{r} = 3(\varphi+\beta)$ and $M_{v} = \varphi+\pi+K_{v}$ are the
SSEE structural constants defined in Section~\ref{sec:axioms}.
One verifies immediately that $1/(2n-1) = -T_{r}/M_{v} = w_{0}$ identically.

The overall normalisation is fixed by requiring that the present-day
dark-energy density equals the SSEE prediction:
%
\begin{equation}
  \rho_{\mathrm{DE},0}
    = \frac{3H_{0}^{2}}{8\pi G}\,\frac{T_{r}}{M_{v}},
  \label{eq:rho_de0}
\end{equation}
%
with $H_{0}$ constrained observationally to
$66.75^{+0.44}_{-0.44}$~km\,s$^{-1}$\,Mpc$^{-1}$ by the MCMC analysis of
Paper~2.  Given the energy scale $\rho_{\mathrm{DE},0}$, the amplitude reads
%
\begin{equation}
  \mathcal{A}
    = \frac{\rho_{\mathrm{DE},0}^{1-n}}{2n-1},
  \label{eq:amplitude}
\end{equation}
%
which completes the specification of $P(X)$ with no free parameters beyond $H_{0}$.

\paragraph{Remark on the CPL evolution.}
The pure k-essence sector~(\ref{eq:w_kessence}) delivers a \emph{constant}
$w = w_{0}$.  The time variation $w_{a}$ arises from the viscous source
(Section~\ref{sec:eft:viscosity}), and is shown there to emerge algebraically
from the scale-factor evolution of the viscosity coefficient.

% ─────────────────────────────────────────────────────────────────────────────
\subsection{Geometric Bulk Viscosity and the Emergence of \texorpdfstring{$w_a$}{wa}}
\label{sec:eft:viscosity}

\subsubsection{Constant-coupling regime ($w_0$)}

The retention constant $\mathrm{KAL}_{0} \equiv \beta + \pi \approx 5.521$
governs the dissipative sector.  Following the Eckart first-order formalism,
the bulk viscosity coefficient is coupled to the instantaneous expansion rate:
%
\begin{equation}
  \zeta_0 = \frac{\mathrm{KAL}_{0}}{8\pi G}\,H.
  \label{eq:zeta}
\end{equation}
%
This delivers a \emph{constant-coupling} viscous pressure
$\Pi_0 = -3\zeta_0 H = -3\,\mathrm{KAL}_{0}H^{2}/8\pi G$,
which governs the background acceleration (Eq.~\ref{eq:Friedmann_a}) and fixes
the present-day effective equation of state $w_0 = -T_r/M_v$ algebraically
(Section~\ref{sec:eft:PX}).

\paragraph{Physical motivation for $\zeta \propto H$.}
Bulk viscosity proportional to the Hubble rate arises naturally in
kinetic theory when the dominant relaxation time is set by the expansion
itself \citep{Weinberg1971, Brevik2005}.  Within SSEE the coupling constant
is not fitted but predicted geometrically as
$\mathrm{KAL}_{0} = (\pi+\varphi)/2 + \pi$.

\subsubsection{Scale-factor evolution and the algebraic determination of \texorpdfstring{$w_a$}{wa}}
\label{sec:eft:wa_derivation}

The constant-coupling viscosity of the preceding section fixes $w_0$
algebraically but predicts a \emph{time-independent} effective equation of
state.  DESI~DR2 data require a time-varying component; within SSEE this must
arise from a scale-factor-dependent extension of the viscosity coefficient.

\paragraph{Structural requirements on $\zeta(a)$.}
Any admissible SSEE viscosity extension $\zeta(a)$ must satisfy three
independent constraints:
\begin{enumerate}[label=(\roman*)]
  \item \textbf{De~Sitter endpoint:} $\zeta(a)\to 0$ as $a\to 1$, so that no
        viscous term shifts the present-day equation of state away from the
        k-essence value $w_0 = -T_r/M_v$.
  \item \textbf{Dimensional consistency:} $[\zeta] = [\rho_{\mathrm{DE}}/H]$,
        constructed from the available dynamical quantities
        $\{\rho_{\mathrm{DE}},H,a\}$.
  \item \textbf{Algebraic closure:} the dimensionless amplitude must be a
        ratio of SSEE structural constants derived exclusively from
        $\{\varphi,\pi\}$.
\end{enumerate}
The unique admissible form satisfying (i) and (ii) at leading order in $(1-a)$
— the natural expansion around the de~Sitter endpoint — is
%
\begin{equation}
  \zeta(a) = \Gamma\,\frac{(1-a)\,\rho_{\mathrm{DE}}(a)}{H(a)},
  \label{eq:zeta_general}
\end{equation}
%
where $\Gamma$ is a dimensionless amplitude to be fixed by constraint (iii).

\paragraph{Algebraic determination of $\Gamma$.}
The k-essence sector has already employed the ratio $T_r/M_v$ to fix
$|w_0|$.  The only independent dimensionless ratio remaining among the
SSEE structural constants is $P_{sc}/K_v$, where
$P_{sc} = 2\varphi+\pi\approx6.378$ (Dynamical Evolution Scalar) and
$K_v = 2(\varphi+\pi)\approx9.519$ (Structural Constraint).
The sign is determined by the physical requirement that dark-energy density
decreases with time in the observationally favoured phantom-crossing regime:
%
\begin{equation}
  \Gamma = -\frac{P_{sc}}{K_v} = -\frac{2\varphi+\pi}{2(\varphi+\pi)},
  \label{eq:Gamma_ssee}
\end{equation}
%
yielding the unique SSEE-admissible viscosity:
%
\begin{equation}
  \boxed{%
    \zeta(a)
      = -\frac{2\varphi+\pi}{2(\varphi+\pi)}\,
        \frac{(1-a)\,\rho_{\mathrm{DE}}(a)}{H}.}
  \label{eq:zeta_ssee}
\end{equation}

\paragraph{Algebraic consistency with CPL.}
The CPL parametrisation $w(a) = w_0 + w_a(1-a)$ is the standard
model-independent description of slowly-varying dark energy; it serves here as
the observational language, not as a physical assumption of SSEE.  Adopting CPL
and substituting $\zeta(a)$ from Eq.~(\ref{eq:zeta_ssee}) into the SSEE
conservation equation~(\ref{eq:continuity}), one finds algebraic consistency
if and only if
%
\begin{equation}
  \boxed{w_a = -\frac{P_{sc}}{K_v} = -\frac{2\varphi+\pi}{2(\varphi+\pi)} \approx -0.6699.}
  \label{eq:wa_ssee}
\end{equation}
%
This is not a derived prediction in the sense that CPL is deduced from the
action; rather, it establishes that \emph{the SSEE viscosity uniquely fixes
the CPL coefficient $w_a$ from the algebraic structure of $\{\varphi,\pi\}$}.
The quantity $w_a$ is therefore not a free parameter of the model but a
computable consequence of the two-axiom system, in the same sense that $w_0$
is fixed by the k-essence exponent $n$.  At $a=1$ the viscous correction
vanishes ($\zeta\to 0$), recovering the pure k-essence regime $w = w_0$.

The implied dark-energy density evolution reads:
%
\begin{equation}
  \rho_{\mathrm{DE}}(a)
    = \rho_{\mathrm{DE},0}\,
      a^{-3(1+w_0+w_a)}\,e^{-3w_a(1-a)},
  \label{eq:rho_CPL}
\end{equation}
%
consistent with the CPL conservation equation
%
\begin{equation}
  \dot\rho_{\mathrm{DE}}
    + 3H\,\rho_{\mathrm{DE}}\,(1+w_0)
    = -3H\,w_a\,(1-a)\,\rho_{\mathrm{DE}},
  \label{eq:cont_CPL}
\end{equation}
%
with the effective viscous pressure
%
\begin{equation}
  \Pi(a)
    = w_a\,(1-a)\,\rho_{\mathrm{DE}}(a).
  \label{eq:Pi_evolving}
\end{equation}
%
At $a \ll 1$ the factor $(1-a)\to 1$ and $\rho_{\mathrm{DE}}$ redshifts
according to Eq.~(\ref{eq:rho_CPL}), so $\zeta$ remains finite and
well-behaved throughout cosmic history.

\paragraph{Boundary conditions.}
The full viscosity interpolates between two algebraically-fixed regimes:
%
\begin{align}
  a \to 0 \;(\text{early}): &\quad
    \zeta \to -\frac{2\varphi+\pi}{2(\varphi+\pi)}\,
              \frac{\rho_{\mathrm{DE,0}}\,a^{-3(1+w_0+w_a)}\,e^{-3w_a}}{H(a)},
  \label{eq:zeta_early}\\[4pt]
  a \to 1 \;(\text{today}): &\quad
    \zeta \to 0.
  \label{eq:zeta_today}
\end{align}
%
Both limits are free of divergences and require no new parameters.

The viscous pressure is therefore
%
\begin{equation}
  \Pi(a)
    = \underbrace{-\frac{3\,\mathrm{KAL}_{0}\,H^{2}}{8\pi G}}_{\text{background}}
      \underbrace{-\;3H\,w_a\,(1-a)\,\rho_{\mathrm{DE}}(a)}_{\text{CPL evolution}},
  \label{eq:Pi_total}
\end{equation}
%
where the first term is the constant-coupling Eckart pressure that sustains
$\ddot a > 0$ and the second term generates the running $w(a)$ observed by
DESI~DR2 \citep{AbdulKarim2025}.

% ─────────────────────────────────────────────────────────────────────────────
\subsection{Modified Friedmann Equations}
\label{sec:eft:friedmann}

Applying the principle of least action ($\delta S = 0$) to the
FLRW metric $ds^{2} = -dt^{2} + a^{2}(t)\,d\mathbf{x}^{2}$ yields
three governing equations for the SSEE background.

\paragraph{I. Energy equation:}
\begin{equation}
  H^{2}
    = \frac{8\pi G}{3}
      \bigl(\rho_{m} + \rho_{r} + \rho_{\mathrm{SSEE}}\bigr).
  \label{eq:Friedmann_H}
\end{equation}

\paragraph{II. Acceleration equation:}
\begin{equation}
  \frac{\ddot{a}}{a}
    = -\frac{4\pi G}{3}
      \!\left[\rho_{m} + 2\rho_{r}
        + \rho_{\mathrm{SSEE}}\!\left(1 + 3w_{0}\right)
        - \frac{9\,\mathrm{KAL}_{0}\,H^{2}}{8\pi G}
      \right].
  \label{eq:Friedmann_a}
\end{equation}

\paragraph{III. Conservation equation with dissipative source:}
\begin{equation}
  \dot{\rho}_{\mathrm{SSEE}}
    + 3H\,\rho_{\mathrm{SSEE}}\!\left(1 + w_{0}\right)
    = \frac{9\,\mathrm{KAL}_{0}\,H^{3}}{8\pi G}.
  \label{eq:continuity}
\end{equation}

\noindent\textbf{Derivation of the coefficient 9.}\par\smallskip
\noindent From Eq.~(\ref{eq:zeta}), the viscous pressure is
$\Pi = -3\mathrm{KAL}_{0}H^{2}/8\pi G$.
In the Raychaudhuri equation the effective pressure of the SSEE sector enters as
$p_{\mathrm{eff}} = w_{0}\rho_{\mathrm{SSEE}} + \Pi$, so the acceleration term becomes
%
\begin{align}
  -\frac{4\pi G}{3}\bigl(\rho_{\mathrm{SSEE}} + 3p_{\mathrm{eff}}\bigr)
  &= -\frac{4\pi G}{3}
     \!\left[\rho_{\mathrm{SSEE}}(1+3w_{0}) + 3\Pi\right] \notag\\
  &= -\frac{4\pi G}{3}
     \!\left[\rho_{\mathrm{SSEE}}(1+3w_{0})
             - \frac{9\,\mathrm{KAL}_{0}\,H^{2}}{8\pi G}\right],
  \label{eq:deriv_coeff9}
\end{align}
which is exactly Eq.~(\ref{eq:Friedmann_a}).  The factor 9 arises from
$3 \times |\Pi|/(H^{2}) = 3 \times 3\mathrm{KAL}_{0}/8\pi G$; a factor
of 3 rather than 9 would correspond to including $\Pi$ only once instead
of as the trace of the viscous stress.

Similarly, the conservation law $\nabla_{\!\mu}T^{\mu\nu}=0$ for the effective
fluid $(w_{0},\Pi)$ gives
$\dot{\rho} + 3H(\rho + w_{0}\rho + \Pi) = 0$, hence
$\dot{\rho} + 3H\rho(1+w_{0}) = 9\mathrm{KAL}_{0}H^{3}/8\pi G$,
confirming that Eq.~(\ref{eq:Friedmann_a}) and Eq.~(\ref{eq:continuity})
are mutually consistent.

% ─────────────────────────────────────────────────────────────────────────────
\subsection{Israel--Stewart Viscous Perturbations}
\label{sec:eft:IS}

The Eckart formulation used in \S\ref{sec:eft:PX}--\ref{sec:eft:friedmann}
is a first-order theory that violates causality \citep{HiscockLindblom1985}:
the viscous signal propagates instantaneously.  Israel--Stewart (IS) theory
\citep{IsraelStewart1979} resolves this by promoting $\Pi$ to a dynamical
variable with a finite relaxation time $\tau_{\Pi}$:
%
\begin{equation}
  \tau_{\Pi} H_{0}\,\tilde{h}(a)\,
    \frac{\mathrm{d}\tilde{\Pi}}{\mathrm{d}\ln a}
  + \tilde{\Pi}
  = -\,\mathrm{KAL}_{0}\,\tilde{h}^{2}(a),
  \label{eq:IS_transport}
\end{equation}
%
where $\tilde{h} = H/H_{0}$, $\tilde{\Pi} = \Pi\,8\pi G/H_{0}^{2}$, and
all densities are normalised to $\rho_{\mathrm{crit},0}$.  In the
limit $\tau_{\Pi} \to 0$, Eq.~(\ref{eq:IS_transport}) reduces to the
Eckart result $\tilde{\Pi} = -\mathrm{KAL}_{0}\tilde{h}^{2}$.

\paragraph{Causality constraint.}
The bulk-viscous sound speed in IS theory is \citep{HiscockLindblom1985}
%
\begin{equation}
  c_{s}^{2}
    = \frac{\zeta}{\rho_{\mathrm{DE}}\,\tau_{\Pi}}
    = \frac{\mathrm{KAL}_{0}}{3\,\Omega_{\mathrm{DE}}\,(\tau_{\Pi} H_{0})}
    \leq 1.
  \label{eq:IS_causality}
\end{equation}
%
Setting $c_{s}^{2} = 1$ (the causality boundary) uniquely fixes the
IS relaxation time as an algebraic combination of SSEE constants:
%
\begin{equation}
  \tau_{\Pi} H_{0}
    \;=\; \frac{\mathrm{KAL}_{0}}{3\,\Omega_{\mathrm{DE}}}
    \;=\; \frac{\mathrm{KAL}_{0}\,M_{v}}{3\,T_{r}}
    \;\approx\; 2.191,
  \label{eq:tau_IS}
\end{equation}
%
with no free parameters.

\paragraph{Growth index.}
Coupling Eq.~(\ref{eq:IS_transport}) to the conservation
equation~(\ref{eq:continuity}) and the linear-growth ODE
%
\begin{equation}
  \delta'' + \left[2 + \frac{\mathrm{d}\ln\tilde{h}}{\mathrm{d}\ln a}\right]\delta'
  = \frac{3}{2}\,\frac{\Omega_{m,\mathrm{dyn}}\,a^{-3}}{\tilde{h}^{2}(a)}\,\delta,
  \label{eq:growth_IS}
\end{equation}
%
we solve the system numerically (\texttt{ssee\_is\_growth.py}) with
the IS background $\tilde{h}^{2}(a) = \Omega_{m,\mathrm{dyn}}\,a^{-3}+\rho_{\mathrm{DE}}(a)$.
The growth rate $f \equiv \mathrm{d}\ln\delta/\mathrm{d}\ln a \approx \Omega_{m}(a)^{\gamma}$
is well described by a power law over $0.1 \leq a \leq 1$ with
%
\begin{equation}
  \gamma_{\mathrm{IS}} = 0.657 \pm 0.002
  \quad(\text{essentially independent of }\tau_{\Pi} H_{0}).
  \label{eq:gamma_IS}
\end{equation}
%
This differs by 6\% from the algebraic sovereignty $\gamma = \varphi^{-1} = 0.618$.
Paper~3 \S5.4 reports both values: the algebraic prediction $\gamma_\mathrm{alg}=0.618$
and the IS numerical result $\gamma_\mathrm{IS}=0.657$.
The IS formalism provides the causal theoretical structure; whether $\gamma = \varphi^{-1}$
is the true attractor of the IS boundary conditions remains a target prediction
for forthcoming Stage-IV weak-lensing surveys (LSST, Euclid).
Note that the direct IS growth-factor integral (this Appendix) yields $S_8=0.837$,
while Paper~3's CAMB post-processing method with the same $\gamma_\mathrm{IS}=0.657$
gives $S_8=0.818$; both are consistent with Planck 2018 ($0.832\pm0.013$) within $1\sigma$.

\paragraph{$S_{8}$ constraint.}
Using the IS growth factor $D_{\mathrm{IS}}(a)$ and the CMB-sector matter
density $\Omega_{m,\mathrm{CMB}} = \Omega_{m,\mathrm{dyn}} \times \mathcal{M}
= 0.3199$ (MIRA, Paper~3),
%
\begin{equation}
  S_{8}^{\mathrm{IS}}
    \equiv \sigma_{8,\mathrm{IS}}
           \sqrt{\frac{\Omega_{m,\mathrm{CMB}}}{0.3}}
    \approx 0.837,
  \label{eq:S8_IS}
\end{equation}
%
within $0.4\sigma$ of the Planck-2018 value $0.832 \pm 0.013$~\citep{PlanckVI2020}.
Using the dynamic sector $\Omega_{m,\mathrm{dyn}} = 0.160$ instead gives
$S_{8} \approx 0.59$ ($6\sigma$ tension); the correct physical choice depends
on the formal resolution of the two-sector Friedmann problem (Section~\ref{sec:eft:status}, item 3).

% ─────────────────────────────────────────────────────────────────────────────
\subsection{Physical Regimes}
\label{sec:eft:regimes}

\paragraph{Late universe ($z \lesssim 2$, DESI~DR2).}
At low redshift $H \approx H_{0}$, so the viscous source
$\propto \mathrm{KAL}_{0}H^{3}/8\pi G$ is small compared to
$\rho_{\mathrm{SSEE}}$.  The acceleration equation is dominated by
$\rho_{\mathrm{SSEE}}(1+3w_{0}) < 0$
(since $w_{0} \approx -0.840 < -1/3$), driving $\ddot{a} > 0$
without dark matter.  This regime is the one tested against DESI~DR2
in Paper~2, where the algebraic prediction $(w_{0},w_{a})$ matches
the data at $0.05\sigma$.

\paragraph{Early universe (CMB epoch, $z \sim 10^{3}$).}
At recombination $H \gg H_{0}$, the geometric retention term
$9\,\mathrm{KAL}_{0}H^{3}/8\pi G$ acts as a dissipative source in
Eq.~(\ref{eq:continuity}), injecting energy into the dark-energy fluid
and modifying its effective equation of state.  This structural
viscosity provides the physical mechanism by which an effective
matter density $\Omega_{m,\mathrm{eff}} = 1 - T_{r}/M_{v} \approx 0.160$
yields the correct CMB acoustic scale through the MIRA factor
$\mathcal{M} = 2 - 1/\varphi^{2} \approx 1.9989$ (Paper~3, \S2.3).

% ─────────────────────────────────────────────────────────────────────────────
% ─────────────────────────────────────────────────────────────────────────────
\subsection{Inflationary Embedding: \texorpdfstring{$\alpha$}{α}-Attractor with \texorpdfstring{$\alpha=\varphi^{4}/3$}{α=φ⁴/3}}
\label{sec:eft:inflation}

The SSEE algebraic predictions for the primordial scalar spectral index
$n_s = 1 - \varphi^{-7}$ (Paper~1, \S\ref{sec:register}) and the
tensor-to-scalar ratio $r = \varphi^{-10}$ (Paper~4) belong simultaneously
to a well-defined class of inflationary models: the
\emph{$\alpha$-attractors} of Kallosh \& Linde \citeyearpar{KalloshLinde2013}.

\paragraph{$\alpha$-attractor predictions.}
In the universal leading-order limit, $\alpha$-attractor models predict
\begin{equation}
  n_s = 1 - \frac{2}{N}, \qquad
  r   = \frac{12\alpha}{N^{2}},
  \label{eq:alpha_attractor}
\end{equation}
where $N$ is the number of inflationary e-folds and $\alpha > 0$
parameterises the inverse curvature of the scalar field's K\"{a}hler target
manifold \citep{KalloshLinde2013,Ferrara2013}.
Setting $\alpha = 1$ recovers Starobinsky $R^{2}$ inflation exactly.

\paragraph{SSEE identification.}
Inserting $n_s = 1 - \varphi^{-7}$ into the first relation gives
\begin{equation}
  N = 2/\varphi^{-7} = 2\varphi^{7} \approx 58.07,
  \label{eq:N_ssee}
\end{equation}
which lies squarely in the observationally required window $50\lesssim N\lesssim 60$.
Substituting $N = 2\varphi^{7}$ and $r = \varphi^{-10}$ into the second relation:
\begin{equation}
  \alpha = \frac{r\,N^{2}}{12}
         = \frac{\varphi^{-10}\cdot 4\varphi^{14}}{12}
         = \frac{4\varphi^{4}}{12}
         = \frac{\varphi^{4}}{3}.
  \label{eq:alpha_ssee}
\end{equation}
This result is algebraically exact (verified numerically to machine precision;
see \texttt{ssee\_inflation\_connection.py}).
Using the Fibonacci identity $\varphi^{4} = 3\varphi + 2$:
\begin{equation}
  \alpha = \frac{3\varphi + 2}{3} = \varphi + \frac{2}{3} \approx 2.285.
  \label{eq:alpha_fibonacci}
\end{equation}

\paragraph{Geometric interpretation.}
In the $\alpha$-attractor framework the K\"{a}hler curvature of the
inflaton's target space (a Poincar\'{e} disk of radius $\sqrt{3\alpha}$)
is $\mathcal{R} = -2/(3\alpha)$.
For $\alpha = \varphi^{4}/3$:
\begin{equation}
  \mathcal{R} = -\frac{2}{3 \cdot \varphi^{4}/3} = -\frac{2}{\varphi^{4}}
               = -\varphi^{-4} \approx -0.146.
  \label{eq:kahler_curvature}
\end{equation}
The inflaton therefore lives on a hyperbolic manifold whose curvature is
set by $\varphi$ alone. This gives a \emph{geometric} origin for the
appearance of $\varphi$ in inflationary observables: $\varphi$ enters not
as a numerical coincidence but as the intrinsic length scale of the
K\"{a}hler geometry of the SSEE inflaton sector.

\paragraph{Slow-roll parameters.}
The SSEE predictions $(n_s,\,r) = (1-\varphi^{-7},\,\varphi^{-10})$ fix the
slow-roll parameters uniquely:
\begin{align}
  \varepsilon &= \frac{r}{16} = \frac{\varphi^{-10}}{16}, \label{eq:sr_eps} \\
  \eta &= \frac{n_s - 1 + 6\varepsilon}{2}
        = \frac{-\varphi^{-7} + 3\varphi^{-10}/8}{2}, \label{eq:sr_eta} \\
  \frac{\eta}{\varepsilon} &= 3 - 8\varphi^{3} = -5 - 16\varphi \approx -30.9.
  \label{eq:sr_ratio}
\end{align}
The ratio $\eta/\varepsilon = 3 - 8\varphi^{3}$ is a pure algebraic
identity derived from the SSEE axioms $\{\varphi, \pi\}$; it defines a
one-parameter family of inflationary potentials $V(\phi_{\rm inf})$ whose
single element consistent with $\alpha$-attractors is $\alpha = \varphi^{4}/3$.

\paragraph{Relation to Starobinsky.}
Starobinsky ($\alpha = 1$) reproduces $n_s$ correctly for $N = 2\varphi^{7}$
but gives $r_{\rm Staro} = 3/\varphi^{14} \approx 0.00356$, a factor of
$\varphi^{4}/3\approx 2.28$ below the SSEE prediction $\varphi^{-10}$.
SSEE is therefore \emph{not} Starobinsky inflation; it is a generalised
$\alpha$-attractor that reduces to Starobinsky only in the limit
$\varphi \to 1$ (i.e.\ $\alpha \to 1$, the fixed point of the golden-ratio
recursion $x = 1 + 1/x$).

% ─────────────────────────────────────────────────────────────────────────────
\subsection{Summary and Open Problems}
\label{sec:eft:status}

Table~\ref{tab:eft_summary} summarises the EFT status.
The action~(\ref{eq:eft_action}) with kinetic exponent~(\ref{eq:n_ssee})
and viscosity~(\ref{eq:zeta_ssee}) constitutes a \emph{complete},
zero-free-parameter EFT description of the SSEE background evolution.
Every constant ($n$, $\mathrm{KAL}_{0}$, $w_0$, $w_a$, $\Omega_{\mathrm{DE}}$) is
derived from $\varphi$ and $\pi$ alone; only the overall energy scale $H_{0}$
enters as an observational input.  Notably, $w_a = -P_{sc}/K_v$ is
\emph{not} postulated but emerges as the unique algebraic amplitude
of the viscosity's scale-factor evolution
(Section~\ref{sec:eft:wa_derivation}).

Four aspects require further development:
%
\begin{enumerate}
  \item \textbf{Growth index.}
        The IS-derived growth index $\gamma_{\mathrm{IS}} = 0.657$
        (\S\ref{sec:eft:IS}) differs by 6\% from the algebraic prediction
        $\varphi^{-1} = 0.618$.  A mechanism that would select
        $\gamma = \varphi^{-1}$ from IS boundary conditions is not yet
        identified; this constitutes a near-prediction testable by
        LSST/Euclid Stage-IV weak lensing.

  \item \textbf{Full perturbation spectrum.}
        The IS growth ODE~(\ref{eq:growth_IS}) governs sub-Hubble dark-matter
        perturbations; the matching to CMB scalar transfer functions and
        the tensor-mode propagation in the IS background require a complete
        Boltzmann implementation.

  \item \textbf{Radiative stability.}
        The EFT cutoff scale $\Lambda_{\mathrm{UV}}$ and the quantum stability
        of the $n < 0$ power-law sector are not yet established.

  \item \textbf{Inflation--dark energy unification.}
        The $\alpha$-attractor sector (\S\ref{sec:eft:inflation}) and the
        k-essence dark-energy sector (\S\ref{sec:eft:PX}) both originate
        from $\varphi$, suggesting a single unified action.
        Constructing the quintessential inflation potential $V(\phi)$
        that reduces to $P(X) = \mathcal{A}X^n$ at late times
        and to an $\alpha=\varphi^4/3$ attractor at early times
        is identified as high-priority future work.
\end{enumerate}

\begin{table}[h]
\centering
\caption{EFT status summary for SSEE-V3.6.}
\label{tab:eft_summary}
\begin{tabular}{lll}
\hline\hline
Element & Status & Derivation \\
\hline
Action form $S$ & \checkmark Complete & Standard GR + k-essence + Eckart \\
$u_{\mu}(\phi)$ coupling & \checkmark Complete & $u_{\mu} = -\partial_{\mu}\phi/\sqrt{2X}$ \\
Kinetic exponent $n$ & \checkmark Algebraic & $n=(T_{r}-M_{v})/2T_{r}$ from $\varphi,\pi$ \\
Equation of state $w_{0}$ & \checkmark Algebraic & $w_{0}=1/(2n-1)=-T_{r}/M_{v}$ \\
Viscosity coupling $\zeta$ & \checkmark Algebraic & $\zeta=\mathrm{KAL}_{0}H/8\pi G$ \\
Friedmann eqs.\ I--III & \checkmark Consistent & Factor 9 verified in \S\ref{sec:eft:friedmann} \\
$w_{a}$ from EFT & \checkmark Algebraic & $\zeta(a)\!\propto\!(1{-}a)\rho_\mathrm{DE}/H$; see \S\ref{sec:eft:wa_derivation} \\
IS relaxation time $\tau_\Pi$ & \checkmark Algebraic & $\tau_\Pi H_0 = \mathrm{KAL}_0/(3\Omega_\mathrm{DE}) \approx 2.191$ \\
IS causality ($c_s^2 \leq 1$) & \checkmark Verified & Eq.~(\ref{eq:IS_causality}); $c_s^2 = 1$ at boundary \\
Growth index $\gamma_\mathrm{IS}$ & \checkmark Numerical & $0.657 \pm 0.002$ (6\% above $\varphi^{-1}=0.618$) \\
$S_8$ (CMB sector) & \checkmark Numerical & $0.837$, within $0.4\sigma$ of Planck \\
$\alpha$-attractor ($\alpha=\varphi^4/3$) & \checkmark Algebraic & Eq.~(\ref{eq:alpha_ssee}); exact, $N=2\varphi^7$ \\
K\"{a}hler curvature $\mathcal{R}=-\varphi^{-4}$ & \checkmark Algebraic & Eq.~(\ref{eq:kahler_curvature}) \\
Full perturbation spectrum & $\circ$ Open & Off-diagonal CMB cov.\ + Stage-IV WL \\
Inflation--DE unification & $\circ$ Open & Quintessential inflation potential \\
Radiative stability & $\circ$ Open & Future work \\
\hline
\end{tabular}
\end{table}

%% All equations verified algebraically (see src/ssee_eft_verify.py)

\section{Effective Field Theory Embedding}
\label{sec:eft}

\textbf{Scope and limitations of this section.}
The background dynamics presented here use the \emph{Eckart first-order}
approximation \citep{Eckart1940}, which is sufficient to derive $w_0$, $w_a$,
and the viscosity coupling $\mathrm{KAL}_0$ on superhorizon scales.
The Eckart formulation is known to violate causality in the relativistic
perturbation sector \citep{HiscockLindblom1985}; for perturbation-level predictions
(B-mode forecasts, sound-speed constraints) the Israel--Stewart (IS) formulation
is required and is developed in Paper~4 \S\ref{sec:is_perturbations}.
The Eckart and IS frameworks agree at leading order in the slow-roll expansion,
so all background results below are unchanged by the IS extension.

\paragraph{Sound speed caveat (EFT embedding).}
For a pure power-law k-essence Lagrangian $P(X)=A\,X^n$ the adiabatic
sound speed satisfies $c_s^2 = 1/(2n-1)$.  With $n\approx-0.095$ (derived below)
this gives $c_s^2\approx-0.84<0$, indicating gradient instability in the
\emph{pure} kinetic sector.  The geometric bulk-viscosity term modifies the
effective perturbation equations and can restore $c_s^2>0$; a full stability
analysis requiring the Israel--Stewart perturbation equations is deferred to
Paper~4.  The EFT embedding presented in this section should be regarded as
a \emph{motivational framework} for the algebraic structure of $w_0$, $w_a$,
not as a complete perturbation-stable theory.

\subsection{The SSEE Action}

The SSEE dark-energy sector is described by a k-essence scalar field $\phi$
minimally coupled to gravity and extended by a geometric bulk-viscosity term
in the Eckart approximation \citep{Eckart1940, Weinberg1971}:
%
\begin{equation}
  S = \int d^{4}x\,\sqrt{-g}
      \left[\frac{R}{16\pi G} + P(X) - \zeta\bigl(\nabla_{\!\mu}u^{\mu}\bigr)^{2}
            + \mathcal{L}_{m}\right],
  \label{eq:eft_action}
\end{equation}
%
where $X \equiv -\tfrac{1}{2}g^{\mu\nu}\partial_{\mu}\phi\,\partial_{\nu}\phi$
is the canonical kinetic term.  The fluid four-velocity of the dark-energy
component is identified with the normalised gradient of the field,
%
\begin{equation}
  u_{\mu} = -\frac{\partial_{\mu}\phi}{\sqrt{2X}},
  \label{eq:four_velocity}
\end{equation}
%
so that $u_{\mu}u^{\mu}=-1$ is satisfied automatically, and the viscous and
kinetic sectors remain self-consistently defined by the single degree of freedom
$\phi$.

% ─────────────────────────────────────────────────────────────────────────────
\subsection{Algebraic Determination of \texorpdfstring{$P(X)$}{P(X)}}
\label{sec:eft:PX}

For a power-law k-essence Lagrangian
$P(X) = \mathcal{A}\,X^{n}$ the equation of state is \emph{exactly} constant:
%
\begin{equation}
  w_{\phi}
    = \frac{P}{2X P_{X} - P}
    = \frac{\mathcal{A}X^{n}}{(2n-1)\mathcal{A}X^{n}}
    = \frac{1}{2n-1}.
  \label{eq:w_kessence}
\end{equation}
%
Imposing the SSEE algebraic constraint $w_{0} = -T_{r}/M_{v}$ uniquely fixes
the kinetic exponent without any fitting:
%
\begin{equation}
  \boxed{%
    n = \frac{T_{r} - M_{v}}{2\,T_{r}}
      = \frac{1 + w_{0}}{2w_{0}}
      \approx -0.09527,}
  \label{eq:n_ssee}
\end{equation}
%
where $T_{r} = 3(\varphi+\beta)$ and $M_{v} = \varphi+\pi+K_{v}$ are the
SSEE structural constants defined in Section~\ref{sec:axioms}.
One verifies immediately that $1/(2n-1) = -T_{r}/M_{v} = w_{0}$ identically.

The overall normalisation is fixed by requiring that the present-day
dark-energy density equals the SSEE prediction:
%
\begin{equation}
  \rho_{\mathrm{DE},0}
    = \frac{3H_{0}^{2}}{8\pi G}\,\frac{T_{r}}{M_{v}},
  \label{eq:rho_de0}
\end{equation}
%
with $H_{0}$ constrained observationally to
$66.75^{+0.44}_{-0.44}$~km\,s$^{-1}$\,Mpc$^{-1}$ by the MCMC analysis of
Paper~2.  Given the energy scale $\rho_{\mathrm{DE},0}$, the amplitude reads
%
\begin{equation}
  \mathcal{A}
    = \frac{\rho_{\mathrm{DE},0}^{1-n}}{2n-1},
  \label{eq:amplitude}
\end{equation}
%
which completes the specification of $P(X)$ with no free parameters beyond $H_{0}$.

\paragraph{Remark on the CPL evolution.}
The pure k-essence sector~(\ref{eq:w_kessence}) delivers a \emph{constant}
$w = w_{0}$.  The time variation $w_{a}$ arises from the viscous source
(Section~\ref{sec:eft:viscosity}), and is shown there to emerge algebraically
from the scale-factor evolution of the viscosity coefficient.

% ─────────────────────────────────────────────────────────────────────────────
\subsection{Geometric Bulk Viscosity and the Emergence of \texorpdfstring{$w_a$}{wa}}
\label{sec:eft:viscosity}

\subsubsection{Constant-coupling regime ($w_0$)}

The retention constant $\mathrm{KAL}_{0} \equiv \beta + \pi \approx 5.521$
governs the dissipative sector.  Following the Eckart first-order formalism,
the bulk viscosity coefficient is coupled to the instantaneous expansion rate:
%
\begin{equation}
  \zeta_0 = \frac{\mathrm{KAL}_{0}}{8\pi G}\,H.
  \label{eq:zeta}
\end{equation}
%
This delivers a \emph{constant-coupling} viscous pressure
$\Pi_0 = -3\zeta_0 H = -3\,\mathrm{KAL}_{0}H^{2}/8\pi G$,
which governs the background acceleration (Eq.~\ref{eq:Friedmann_a}) and fixes
the present-day effective equation of state $w_0 = -T_r/M_v$ algebraically
(Section~\ref{sec:eft:PX}).

\paragraph{Physical motivation for $\zeta \propto H$.}
Bulk viscosity proportional to the Hubble rate arises naturally in
kinetic theory when the dominant relaxation time is set by the expansion
itself \citep{Weinberg1971, Brevik2005}.  Within SSEE the coupling constant
is not fitted but predicted geometrically as
$\mathrm{KAL}_{0} = (\pi+\varphi)/2 + \pi$.

\subsubsection{Scale-factor evolution and the algebraic determination of \texorpdfstring{$w_a$}{wa}}
\label{sec:eft:wa_derivation}

The constant-coupling viscosity of the preceding section fixes $w_0$
algebraically but predicts a \emph{time-independent} effective equation of
state.  DESI~DR2 data require a time-varying component; within SSEE this must
arise from a scale-factor-dependent extension of the viscosity coefficient.

\paragraph{Structural requirements on $\zeta(a)$.}
Any admissible SSEE viscosity extension $\zeta(a)$ must satisfy three
independent constraints:
\begin{enumerate}[label=(\roman*)]
  \item \textbf{De~Sitter endpoint:} $\zeta(a)\to 0$ as $a\to 1$, so that no
        viscous term shifts the present-day equation of state away from the
        k-essence value $w_0 = -T_r/M_v$.
  \item \textbf{Dimensional consistency:} $[\zeta] = [\rho_{\mathrm{DE}}/H]$,
        constructed from the available dynamical quantities
        $\{\rho_{\mathrm{DE}},H,a\}$.
  \item \textbf{Algebraic closure:} the dimensionless amplitude must be a
        ratio of SSEE structural constants derived exclusively from
        $\{\varphi,\pi\}$.
\end{enumerate}
The unique admissible form satisfying (i) and (ii) at leading order in $(1-a)$
— the natural expansion around the de~Sitter endpoint — is
%
\begin{equation}
  \zeta(a) = \Gamma\,\frac{(1-a)\,\rho_{\mathrm{DE}}(a)}{H(a)},
  \label{eq:zeta_general}
\end{equation}
%
where $\Gamma$ is a dimensionless amplitude to be fixed by constraint (iii).

\paragraph{Algebraic determination of $\Gamma$.}
The k-essence sector has already employed the ratio $T_r/M_v$ to fix
$|w_0|$.  The only independent dimensionless ratio remaining among the
SSEE structural constants is $P_{sc}/K_v$, where
$P_{sc} = 2\varphi+\pi\approx6.378$ (Dynamical Evolution Scalar) and
$K_v = 2(\varphi+\pi)\approx9.519$ (Structural Constraint).
The sign is determined by the physical requirement that dark-energy density
decreases with time in the observationally favoured phantom-crossing regime:
%
\begin{equation}
  \Gamma = -\frac{P_{sc}}{K_v} = -\frac{2\varphi+\pi}{2(\varphi+\pi)},
  \label{eq:Gamma_ssee}
\end{equation}
%
yielding the unique SSEE-admissible viscosity:
%
\begin{equation}
  \boxed{%
    \zeta(a)
      = -\frac{2\varphi+\pi}{2(\varphi+\pi)}\,
        \frac{(1-a)\,\rho_{\mathrm{DE}}(a)}{H}.}
  \label{eq:zeta_ssee}
\end{equation}

\paragraph{Algebraic consistency with CPL.}
The CPL parametrisation $w(a) = w_0 + w_a(1-a)$ is the standard
model-independent description of slowly-varying dark energy; it serves here as
the observational language, not as a physical assumption of SSEE.  Adopting CPL
and substituting $\zeta(a)$ from Eq.~(\ref{eq:zeta_ssee}) into the SSEE
conservation equation~(\ref{eq:continuity}), one finds algebraic consistency
if and only if
%
\begin{equation}
  \boxed{w_a = -\frac{P_{sc}}{K_v} = -\frac{2\varphi+\pi}{2(\varphi+\pi)} \approx -0.6699.}
  \label{eq:wa_ssee}
\end{equation}
%
This is not a derived prediction in the sense that CPL is deduced from the
action; rather, it establishes that \emph{the SSEE viscosity uniquely fixes
the CPL coefficient $w_a$ from the algebraic structure of $\{\varphi,\pi\}$}.
The quantity $w_a$ is therefore not a free parameter of the model but a
computable consequence of the two-axiom system, in the same sense that $w_0$
is fixed by the k-essence exponent $n$.  At $a=1$ the viscous correction
vanishes ($\zeta\to 0$), recovering the pure k-essence regime $w = w_0$.

The implied dark-energy density evolution reads:
%
\begin{equation}
  \rho_{\mathrm{DE}}(a)
    = \rho_{\mathrm{DE},0}\,
      a^{-3(1+w_0+w_a)}\,e^{-3w_a(1-a)},
  \label{eq:rho_CPL}
\end{equation}
%
consistent with the CPL conservation equation
%
\begin{equation}
  \dot\rho_{\mathrm{DE}}
    + 3H\,\rho_{\mathrm{DE}}\,(1+w_0)
    = -3H\,w_a\,(1-a)\,\rho_{\mathrm{DE}},
  \label{eq:cont_CPL}
\end{equation}
%
with the effective viscous pressure
%
\begin{equation}
  \Pi(a)
    = w_a\,(1-a)\,\rho_{\mathrm{DE}}(a).
  \label{eq:Pi_evolving}
\end{equation}
%
At $a \ll 1$ the factor $(1-a)\to 1$ and $\rho_{\mathrm{DE}}$ redshifts
according to Eq.~(\ref{eq:rho_CPL}), so $\zeta$ remains finite and
well-behaved throughout cosmic history.

\paragraph{Boundary conditions.}
The full viscosity interpolates between two algebraically-fixed regimes:
%
\begin{align}
  a \to 0 \;(\text{early}): &\quad
    \zeta \to -\frac{2\varphi+\pi}{2(\varphi+\pi)}\,
              \frac{\rho_{\mathrm{DE,0}}\,a^{-3(1+w_0+w_a)}\,e^{-3w_a}}{H(a)},
  \label{eq:zeta_early}\\[4pt]
  a \to 1 \;(\text{today}): &\quad
    \zeta \to 0.
  \label{eq:zeta_today}
\end{align}
%
Both limits are free of divergences and require no new parameters.

The viscous pressure is therefore
%
\begin{equation}
  \Pi(a)
    = \underbrace{-\frac{3\,\mathrm{KAL}_{0}\,H^{2}}{8\pi G}}_{\text{background}}
      \underbrace{-\;3H\,w_a\,(1-a)\,\rho_{\mathrm{DE}}(a)}_{\text{CPL evolution}},
  \label{eq:Pi_total}
\end{equation}
%
where the first term is the constant-coupling Eckart pressure that sustains
$\ddot a > 0$ and the second term generates the running $w(a)$ observed by
DESI~DR2 \citep{AbdulKarim2025}.

% ─────────────────────────────────────────────────────────────────────────────
\subsection{Modified Friedmann Equations}
\label{sec:eft:friedmann}

Applying the principle of least action ($\delta S = 0$) to the
FLRW metric $ds^{2} = -dt^{2} + a^{2}(t)\,d\mathbf{x}^{2}$ yields
three governing equations for the SSEE background.

\paragraph{I. Energy equation:}
\begin{equation}
  H^{2}
    = \frac{8\pi G}{3}
      \bigl(\rho_{m} + \rho_{r} + \rho_{\mathrm{SSEE}}\bigr).
  \label{eq:Friedmann_H}
\end{equation}

\paragraph{II. Acceleration equation:}
\begin{equation}
  \frac{\ddot{a}}{a}
    = -\frac{4\pi G}{3}
      \!\left[\rho_{m} + 2\rho_{r}
        + \rho_{\mathrm{SSEE}}\!\left(1 + 3w_{0}\right)
        - \frac{9\,\mathrm{KAL}_{0}\,H^{2}}{8\pi G}
      \right].
  \label{eq:Friedmann_a}
\end{equation}

\paragraph{III. Conservation equation with dissipative source:}
\begin{equation}
  \dot{\rho}_{\mathrm{SSEE}}
    + 3H\,\rho_{\mathrm{SSEE}}\!\left(1 + w_{0}\right)
    = \frac{9\,\mathrm{KAL}_{0}\,H^{3}}{8\pi G}.
  \label{eq:continuity}
\end{equation}

\noindent\textbf{Derivation of the coefficient 9.}\par\smallskip
\noindent From Eq.~(\ref{eq:zeta}), the viscous pressure is
$\Pi = -3\mathrm{KAL}_{0}H^{2}/8\pi G$.
In the Raychaudhuri equation the effective pressure of the SSEE sector enters as
$p_{\mathrm{eff}} = w_{0}\rho_{\mathrm{SSEE}} + \Pi$, so the acceleration term becomes
%
\begin{align}
  -\frac{4\pi G}{3}\bigl(\rho_{\mathrm{SSEE}} + 3p_{\mathrm{eff}}\bigr)
  &= -\frac{4\pi G}{3}
     \!\left[\rho_{\mathrm{SSEE}}(1+3w_{0}) + 3\Pi\right] \notag\\
  &= -\frac{4\pi G}{3}
     \!\left[\rho_{\mathrm{SSEE}}(1+3w_{0})
             - \frac{9\,\mathrm{KAL}_{0}\,H^{2}}{8\pi G}\right],
  \label{eq:deriv_coeff9}
\end{align}
which is exactly Eq.~(\ref{eq:Friedmann_a}).  The factor 9 arises from
$3 \times |\Pi|/(H^{2}) = 3 \times 3\mathrm{KAL}_{0}/8\pi G$; a factor
of 3 rather than 9 would correspond to including $\Pi$ only once instead
of as the trace of the viscous stress.

Similarly, the conservation law $\nabla_{\!\mu}T^{\mu\nu}=0$ for the effective
fluid $(w_{0},\Pi)$ gives
$\dot{\rho} + 3H(\rho + w_{0}\rho + \Pi) = 0$, hence
$\dot{\rho} + 3H\rho(1+w_{0}) = 9\mathrm{KAL}_{0}H^{3}/8\pi G$,
confirming that Eq.~(\ref{eq:Friedmann_a}) and Eq.~(\ref{eq:continuity})
are mutually consistent.

% ─────────────────────────────────────────────────────────────────────────────
\subsection{Israel--Stewart Viscous Perturbations}
\label{sec:eft:IS}

The Eckart formulation used in \S\ref{sec:eft:PX}--\ref{sec:eft:friedmann}
is a first-order theory that violates causality \citep{HiscockLindblom1985}:
the viscous signal propagates instantaneously.  Israel--Stewart (IS) theory
\citep{IsraelStewart1979} resolves this by promoting $\Pi$ to a dynamical
variable with a finite relaxation time $\tau_{\Pi}$:
%
\begin{equation}
  \tau_{\Pi} H_{0}\,\tilde{h}(a)\,
    \frac{\mathrm{d}\tilde{\Pi}}{\mathrm{d}\ln a}
  + \tilde{\Pi}
  = -\,\mathrm{KAL}_{0}\,\tilde{h}^{2}(a),
  \label{eq:IS_transport}
\end{equation}
%
where $\tilde{h} = H/H_{0}$, $\tilde{\Pi} = \Pi\,8\pi G/H_{0}^{2}$, and
all densities are normalised to $\rho_{\mathrm{crit},0}$.  In the
limit $\tau_{\Pi} \to 0$, Eq.~(\ref{eq:IS_transport}) reduces to the
Eckart result $\tilde{\Pi} = -\mathrm{KAL}_{0}\tilde{h}^{2}$.

\paragraph{Causality constraint.}
The bulk-viscous sound speed in IS theory is \citep{HiscockLindblom1985}
%
\begin{equation}
  c_{s}^{2}
    = \frac{\zeta}{\rho_{\mathrm{DE}}\,\tau_{\Pi}}
    = \frac{\mathrm{KAL}_{0}}{3\,\Omega_{\mathrm{DE}}\,(\tau_{\Pi} H_{0})}
    \leq 1.
  \label{eq:IS_causality}
\end{equation}
%
Setting $c_{s}^{2} = 1$ (the causality boundary) uniquely fixes the
IS relaxation time as an algebraic combination of SSEE constants:
%
\begin{equation}
  \tau_{\Pi} H_{0}
    \;=\; \frac{\mathrm{KAL}_{0}}{3\,\Omega_{\mathrm{DE}}}
    \;=\; \frac{\mathrm{KAL}_{0}\,M_{v}}{3\,T_{r}}
    \;\approx\; 2.191,
  \label{eq:tau_IS}
\end{equation}
%
with no free parameters.

\paragraph{Growth index.}
Coupling Eq.~(\ref{eq:IS_transport}) to the conservation
equation~(\ref{eq:continuity}) and the linear-growth ODE
%
\begin{equation}
  \delta'' + \left[2 + \frac{\mathrm{d}\ln\tilde{h}}{\mathrm{d}\ln a}\right]\delta'
  = \frac{3}{2}\,\frac{\Omega_{m,\mathrm{dyn}}\,a^{-3}}{\tilde{h}^{2}(a)}\,\delta,
  \label{eq:growth_IS}
\end{equation}
%
we solve the system numerically (\texttt{ssee\_is\_growth.py}) with
the IS background $\tilde{h}^{2}(a) = \Omega_{m,\mathrm{dyn}}\,a^{-3}+\rho_{\mathrm{DE}}(a)$.
The growth rate $f \equiv \mathrm{d}\ln\delta/\mathrm{d}\ln a \approx \Omega_{m}(a)^{\gamma}$
is well described by a power law over $0.1 \leq a \leq 1$ with
%
\begin{equation}
  \gamma_{\mathrm{IS}} = 0.657 \pm 0.002
  \quad(\text{essentially independent of }\tau_{\Pi} H_{0}).
  \label{eq:gamma_IS}
\end{equation}
%
This differs by 6\% from the algebraic sovereignty $\gamma = \varphi^{-1} = 0.618$.
Paper~3 \S5.4 reports both values: the algebraic prediction $\gamma_\mathrm{alg}=0.618$
and the IS numerical result $\gamma_\mathrm{IS}=0.657$.
The IS formalism provides the causal theoretical structure; whether $\gamma = \varphi^{-1}$
is the true attractor of the IS boundary conditions remains a target prediction
for forthcoming Stage-IV weak-lensing surveys (LSST, Euclid).
Note that the direct IS growth-factor integral (this Appendix) yields $S_8=0.837$,
while Paper~3's CAMB post-processing method with the same $\gamma_\mathrm{IS}=0.657$
gives $S_8=0.818$; both are consistent with Planck 2018 ($0.832\pm0.013$) within $1\sigma$.

\paragraph{$S_{8}$ constraint.}
Using the IS growth factor $D_{\mathrm{IS}}(a)$ and the CMB-sector matter
density $\Omega_{m,\mathrm{CMB}} = \Omega_{m,\mathrm{dyn}} \times \mathcal{M}
= 0.3199$ (MIRA, Paper~3),
%
\begin{equation}
  S_{8}^{\mathrm{IS}}
    \equiv \sigma_{8,\mathrm{IS}}
           \sqrt{\frac{\Omega_{m,\mathrm{CMB}}}{0.3}}
    \approx 0.837,
  \label{eq:S8_IS}
\end{equation}
%
within $0.4\sigma$ of the Planck-2018 value $0.832 \pm 0.013$~\citep{PlanckVI2020}.
Using the dynamic sector $\Omega_{m,\mathrm{dyn}} = 0.160$ instead gives
$S_{8} \approx 0.59$ ($6\sigma$ tension); the correct physical choice depends
on the formal resolution of the two-sector Friedmann problem (Section~\ref{sec:eft:status}, item 3).

% ─────────────────────────────────────────────────────────────────────────────
\subsection{Physical Regimes}
\label{sec:eft:regimes}

\paragraph{Late universe ($z \lesssim 2$, DESI~DR2).}
At low redshift $H \approx H_{0}$, so the viscous source
$\propto \mathrm{KAL}_{0}H^{3}/8\pi G$ is small compared to
$\rho_{\mathrm{SSEE}}$.  The acceleration equation is dominated by
$\rho_{\mathrm{SSEE}}(1+3w_{0}) < 0$
(since $w_{0} \approx -0.840 < -1/3$), driving $\ddot{a} > 0$
without dark matter.  This regime is the one tested against DESI~DR2
in Paper~2, where the algebraic prediction $(w_{0},w_{a})$ matches
the data at $0.05\sigma$.

\paragraph{Early universe (CMB epoch, $z \sim 10^{3}$).}
At recombination $H \gg H_{0}$, the geometric retention term
$9\,\mathrm{KAL}_{0}H^{3}/8\pi G$ acts as a dissipative source in
Eq.~(\ref{eq:continuity}), injecting energy into the dark-energy fluid
and modifying its effective equation of state.  This structural
viscosity provides the physical mechanism by which an effective
matter density $\Omega_{m,\mathrm{eff}} = 1 - T_{r}/M_{v} \approx 0.160$
yields the correct CMB acoustic scale through the MIRA factor
$\mathcal{M} = 2 - 1/\varphi^{2} \approx 1.9989$ (Paper~3, \S2.3).

% ─────────────────────────────────────────────────────────────────────────────
% ─────────────────────────────────────────────────────────────────────────────
\subsection{Inflationary Embedding: \texorpdfstring{$\alpha$}{α}-Attractor with \texorpdfstring{$\alpha=\varphi^{4}/3$}{α=φ⁴/3}}
\label{sec:eft:inflation}

The SSEE algebraic predictions for the primordial scalar spectral index
$n_s = 1 - \varphi^{-7}$ (Paper~1, \S\ref{sec:register}) and the
tensor-to-scalar ratio $r = \varphi^{-10}$ (Paper~4) belong simultaneously
to a well-defined class of inflationary models: the
\emph{$\alpha$-attractors} of Kallosh \& Linde \citeyearpar{KalloshLinde2013}.

\paragraph{$\alpha$-attractor predictions.}
In the universal leading-order limit, $\alpha$-attractor models predict
\begin{equation}
  n_s = 1 - \frac{2}{N}, \qquad
  r   = \frac{12\alpha}{N^{2}},
  \label{eq:alpha_attractor}
\end{equation}
where $N$ is the number of inflationary e-folds and $\alpha > 0$
parameterises the inverse curvature of the scalar field's K\"{a}hler target
manifold \citep{KalloshLinde2013,Ferrara2013}.
Setting $\alpha = 1$ recovers Starobinsky $R^{2}$ inflation exactly.

\paragraph{SSEE identification.}
Inserting $n_s = 1 - \varphi^{-7}$ into the first relation gives
\begin{equation}
  N = 2/\varphi^{-7} = 2\varphi^{7} \approx 58.07,
  \label{eq:N_ssee}
\end{equation}
which lies squarely in the observationally required window $50\lesssim N\lesssim 60$.
Substituting $N = 2\varphi^{7}$ and $r = \varphi^{-10}$ into the second relation:
\begin{equation}
  \alpha = \frac{r\,N^{2}}{12}
         = \frac{\varphi^{-10}\cdot 4\varphi^{14}}{12}
         = \frac{4\varphi^{4}}{12}
         = \frac{\varphi^{4}}{3}.
  \label{eq:alpha_ssee}
\end{equation}
This result is algebraically exact (verified numerically to machine precision;
see \texttt{ssee\_inflation\_connection.py}).
Using the Fibonacci identity $\varphi^{4} = 3\varphi + 2$:
\begin{equation}
  \alpha = \frac{3\varphi + 2}{3} = \varphi + \frac{2}{3} \approx 2.285.
  \label{eq:alpha_fibonacci}
\end{equation}

\paragraph{Geometric interpretation.}
In the $\alpha$-attractor framework the K\"{a}hler curvature of the
inflaton's target space (a Poincar\'{e} disk of radius $\sqrt{3\alpha}$)
is $\mathcal{R} = -2/(3\alpha)$.
For $\alpha = \varphi^{4}/3$:
\begin{equation}
  \mathcal{R} = -\frac{2}{3 \cdot \varphi^{4}/3} = -\frac{2}{\varphi^{4}}
               = -\varphi^{-4} \approx -0.146.
  \label{eq:kahler_curvature}
\end{equation}
The inflaton therefore lives on a hyperbolic manifold whose curvature is
set by $\varphi$ alone. This gives a \emph{geometric} origin for the
appearance of $\varphi$ in inflationary observables: $\varphi$ enters not
as a numerical coincidence but as the intrinsic length scale of the
K\"{a}hler geometry of the SSEE inflaton sector.

\paragraph{Slow-roll parameters.}
The SSEE predictions $(n_s,\,r) = (1-\varphi^{-7},\,\varphi^{-10})$ fix the
slow-roll parameters uniquely:
\begin{align}
  \varepsilon &= \frac{r}{16} = \frac{\varphi^{-10}}{16}, \label{eq:sr_eps} \\
  \eta &= \frac{n_s - 1 + 6\varepsilon}{2}
        = \frac{-\varphi^{-7} + 3\varphi^{-10}/8}{2}, \label{eq:sr_eta} \\
  \frac{\eta}{\varepsilon} &= 3 - 8\varphi^{3} = -5 - 16\varphi \approx -30.9.
  \label{eq:sr_ratio}
\end{align}
The ratio $\eta/\varepsilon = 3 - 8\varphi^{3}$ is a pure algebraic
identity derived from the SSEE axioms $\{\varphi, \pi\}$; it defines a
one-parameter family of inflationary potentials $V(\phi_{\rm inf})$ whose
single element consistent with $\alpha$-attractors is $\alpha = \varphi^{4}/3$.

\paragraph{Relation to Starobinsky.}
Starobinsky ($\alpha = 1$) reproduces $n_s$ correctly for $N = 2\varphi^{7}$
but gives $r_{\rm Staro} = 3/\varphi^{14} \approx 0.00356$, a factor of
$\varphi^{4}/3\approx 2.28$ below the SSEE prediction $\varphi^{-10}$.
SSEE is therefore \emph{not} Starobinsky inflation; it is a generalised
$\alpha$-attractor that reduces to Starobinsky only in the limit
$\varphi \to 1$ (i.e.\ $\alpha \to 1$, the fixed point of the golden-ratio
recursion $x = 1 + 1/x$).

% ─────────────────────────────────────────────────────────────────────────────
\subsection{Summary and Open Problems}
\label{sec:eft:status}

Table~\ref{tab:eft_summary} summarises the EFT status.
The action~(\ref{eq:eft_action}) with kinetic exponent~(\ref{eq:n_ssee})
and viscosity~(\ref{eq:zeta_ssee}) constitutes a \emph{complete},
zero-free-parameter EFT description of the SSEE background evolution.
Every constant ($n$, $\mathrm{KAL}_{0}$, $w_0$, $w_a$, $\Omega_{\mathrm{DE}}$) is
derived from $\varphi$ and $\pi$ alone; only the overall energy scale $H_{0}$
enters as an observational input.  Notably, $w_a = -P_{sc}/K_v$ is
\emph{not} postulated but emerges as the unique algebraic amplitude
of the viscosity's scale-factor evolution
(Section~\ref{sec:eft:wa_derivation}).

Four aspects require further development:
%
\begin{enumerate}
  \item \textbf{Growth index.}
        The IS-derived growth index $\gamma_{\mathrm{IS}} = 0.657$
        (\S\ref{sec:eft:IS}) differs by 6\% from the algebraic prediction
        $\varphi^{-1} = 0.618$.  A mechanism that would select
        $\gamma = \varphi^{-1}$ from IS boundary conditions is not yet
        identified; this constitutes a near-prediction testable by
        LSST/Euclid Stage-IV weak lensing.

  \item \textbf{Full perturbation spectrum.}
        The IS growth ODE~(\ref{eq:growth_IS}) governs sub-Hubble dark-matter
        perturbations; the matching to CMB scalar transfer functions and
        the tensor-mode propagation in the IS background require a complete
        Boltzmann implementation.

  \item \textbf{Radiative stability.}
        The EFT cutoff scale $\Lambda_{\mathrm{UV}}$ and the quantum stability
        of the $n < 0$ power-law sector are not yet established.

  \item \textbf{Inflation--dark energy unification.}
        The $\alpha$-attractor sector (\S\ref{sec:eft:inflation}) and the
        k-essence dark-energy sector (\S\ref{sec:eft:PX}) both originate
        from $\varphi$, suggesting a single unified action.
        Constructing the quintessential inflation potential $V(\phi)$
        that reduces to $P(X) = \mathcal{A}X^n$ at late times
        and to an $\alpha=\varphi^4/3$ attractor at early times
        is identified as high-priority future work.
\end{enumerate}

\begin{table}[h]
\centering
\caption{EFT status summary for SSEE-V3.6.}
\label{tab:eft_summary}
\begin{tabular}{lll}
\hline\hline
Element & Status & Derivation \\
\hline
Action form $S$ & \checkmark Complete & Standard GR + k-essence + Eckart \\
$u_{\mu}(\phi)$ coupling & \checkmark Complete & $u_{\mu} = -\partial_{\mu}\phi/\sqrt{2X}$ \\
Kinetic exponent $n$ & \checkmark Algebraic & $n=(T_{r}-M_{v})/2T_{r}$ from $\varphi,\pi$ \\
Equation of state $w_{0}$ & \checkmark Algebraic & $w_{0}=1/(2n-1)=-T_{r}/M_{v}$ \\
Viscosity coupling $\zeta$ & \checkmark Algebraic & $\zeta=\mathrm{KAL}_{0}H/8\pi G$ \\
Friedmann eqs.\ I--III & \checkmark Consistent & Factor 9 verified in \S\ref{sec:eft:friedmann} \\
$w_{a}$ from EFT & \checkmark Algebraic & $\zeta(a)\!\propto\!(1{-}a)\rho_\mathrm{DE}/H$; see \S\ref{sec:eft:wa_derivation} \\
IS relaxation time $\tau_\Pi$ & \checkmark Algebraic & $\tau_\Pi H_0 = \mathrm{KAL}_0/(3\Omega_\mathrm{DE}) \approx 2.191$ \\
IS causality ($c_s^2 \leq 1$) & \checkmark Verified & Eq.~(\ref{eq:IS_causality}); $c_s^2 = 1$ at boundary \\
Growth index $\gamma_\mathrm{IS}$ & \checkmark Numerical & $0.657 \pm 0.002$ (6\% above $\varphi^{-1}=0.618$) \\
$S_8$ (CMB sector) & \checkmark Numerical & $0.837$, within $0.4\sigma$ of Planck \\
$\alpha$-attractor ($\alpha=\varphi^4/3$) & \checkmark Algebraic & Eq.~(\ref{eq:alpha_ssee}); exact, $N=2\varphi^7$ \\
K\"{a}hler curvature $\mathcal{R}=-\varphi^{-4}$ & \checkmark Algebraic & Eq.~(\ref{eq:kahler_curvature}) \\
Full perturbation spectrum & $\circ$ Open & Off-diagonal CMB cov.\ + Stage-IV WL \\
Inflation--DE unification & $\circ$ Open & Quintessential inflation potential \\
Radiative stability & $\circ$ Open & Future work \\
\hline
\end{tabular}
\end{table}

%% All equations verified algebraically (see src/ssee_eft_verify.py)

\section{Effective Field Theory Embedding}
\label{sec:eft}

\textbf{Scope and limitations of this section.}
The background dynamics presented here use the \emph{Eckart first-order}
approximation \citep{Eckart1940}, which is sufficient to derive $w_0$, $w_a$,
and the viscosity coupling $\mathrm{KAL}_0$ on superhorizon scales.
The Eckart formulation is known to violate causality in the relativistic
perturbation sector \citep{HiscockLindblom1985}; for perturbation-level predictions
(B-mode forecasts, sound-speed constraints) the Israel--Stewart (IS) formulation
is required and is developed in Paper~4 \S\ref{sec:is_perturbations}.
The Eckart and IS frameworks agree at leading order in the slow-roll expansion,
so all background results below are unchanged by the IS extension.

\paragraph{Sound speed caveat (EFT embedding).}
For a pure power-law k-essence Lagrangian $P(X)=A\,X^n$ the adiabatic
sound speed satisfies $c_s^2 = 1/(2n-1)$.  With $n\approx-0.095$ (derived below)
this gives $c_s^2\approx-0.84<0$, indicating gradient instability in the
\emph{pure} kinetic sector.  The geometric bulk-viscosity term modifies the
effective perturbation equations and can restore $c_s^2>0$; a full stability
analysis requiring the Israel--Stewart perturbation equations is deferred to
Paper~4.  The EFT embedding presented in this section should be regarded as
a \emph{motivational framework} for the algebraic structure of $w_0$, $w_a$,
not as a complete perturbation-stable theory.

\subsection{The SSEE Action}

The SSEE dark-energy sector is described by a k-essence scalar field $\phi$
minimally coupled to gravity and extended by a geometric bulk-viscosity term
in the Eckart approximation \citep{Eckart1940, Weinberg1971}:
%
\begin{equation}
  S = \int d^{4}x\,\sqrt{-g}
      \left[\frac{R}{16\pi G} + P(X) - \zeta\bigl(\nabla_{\!\mu}u^{\mu}\bigr)^{2}
            + \mathcal{L}_{m}\right],
  \label{eq:eft_action}
\end{equation}
%
where $X \equiv -\tfrac{1}{2}g^{\mu\nu}\partial_{\mu}\phi\,\partial_{\nu}\phi$
is the canonical kinetic term.  The fluid four-velocity of the dark-energy
component is identified with the normalised gradient of the field,
%
\begin{equation}
  u_{\mu} = -\frac{\partial_{\mu}\phi}{\sqrt{2X}},
  \label{eq:four_velocity}
\end{equation}
%
so that $u_{\mu}u^{\mu}=-1$ is satisfied automatically, and the viscous and
kinetic sectors remain self-consistently defined by the single degree of freedom
$\phi$.

% ─────────────────────────────────────────────────────────────────────────────
\subsection{Algebraic Determination of \texorpdfstring{$P(X)$}{P(X)}}
\label{sec:eft:PX}

For a power-law k-essence Lagrangian
$P(X) = \mathcal{A}\,X^{n}$ the equation of state is \emph{exactly} constant:
%
\begin{equation}
  w_{\phi}
    = \frac{P}{2X P_{X} - P}
    = \frac{\mathcal{A}X^{n}}{(2n-1)\mathcal{A}X^{n}}
    = \frac{1}{2n-1}.
  \label{eq:w_kessence}
\end{equation}
%
Imposing the SSEE algebraic constraint $w_{0} = -T_{r}/M_{v}$ uniquely fixes
the kinetic exponent without any fitting:
%
\begin{equation}
  \boxed{%
    n = \frac{T_{r} - M_{v}}{2\,T_{r}}
      = \frac{1 + w_{0}}{2w_{0}}
      \approx -0.09527,}
  \label{eq:n_ssee}
\end{equation}
%
where $T_{r} = 3(\varphi+\beta)$ and $M_{v} = \varphi+\pi+K_{v}$ are the
SSEE structural constants defined in Section~\ref{sec:axioms}.
One verifies immediately that $1/(2n-1) = -T_{r}/M_{v} = w_{0}$ identically.

The overall normalisation is fixed by requiring that the present-day
dark-energy density equals the SSEE prediction:
%
\begin{equation}
  \rho_{\mathrm{DE},0}
    = \frac{3H_{0}^{2}}{8\pi G}\,\frac{T_{r}}{M_{v}},
  \label{eq:rho_de0}
\end{equation}
%
with $H_{0}$ constrained observationally to
$66.75^{+0.44}_{-0.44}$~km\,s$^{-1}$\,Mpc$^{-1}$ by the MCMC analysis of
Paper~2.  Given the energy scale $\rho_{\mathrm{DE},0}$, the amplitude reads
%
\begin{equation}
  \mathcal{A}
    = \frac{\rho_{\mathrm{DE},0}^{1-n}}{2n-1},
  \label{eq:amplitude}
\end{equation}
%
which completes the specification of $P(X)$ with no free parameters beyond $H_{0}$.

\paragraph{Remark on the CPL evolution.}
The pure k-essence sector~(\ref{eq:w_kessence}) delivers a \emph{constant}
$w = w_{0}$.  The time variation $w_{a}$ arises from the viscous source
(Section~\ref{sec:eft:viscosity}), and is shown there to emerge algebraically
from the scale-factor evolution of the viscosity coefficient.

% ─────────────────────────────────────────────────────────────────────────────
\subsection{Geometric Bulk Viscosity and the Emergence of \texorpdfstring{$w_a$}{wa}}
\label{sec:eft:viscosity}

\subsubsection{Constant-coupling regime ($w_0$)}

The retention constant $\mathrm{KAL}_{0} \equiv \beta + \pi \approx 5.521$
governs the dissipative sector.  Following the Eckart first-order formalism,
the bulk viscosity coefficient is coupled to the instantaneous expansion rate:
%
\begin{equation}
  \zeta_0 = \frac{\mathrm{KAL}_{0}}{8\pi G}\,H.
  \label{eq:zeta}
\end{equation}
%
This delivers a \emph{constant-coupling} viscous pressure
$\Pi_0 = -3\zeta_0 H = -3\,\mathrm{KAL}_{0}H^{2}/8\pi G$,
which governs the background acceleration (Eq.~\ref{eq:Friedmann_a}) and fixes
the present-day effective equation of state $w_0 = -T_r/M_v$ algebraically
(Section~\ref{sec:eft:PX}).

\paragraph{Physical motivation for $\zeta \propto H$.}
Bulk viscosity proportional to the Hubble rate arises naturally in
kinetic theory when the dominant relaxation time is set by the expansion
itself \citep{Weinberg1971, Brevik2005}.  Within SSEE the coupling constant
is not fitted but predicted geometrically as
$\mathrm{KAL}_{0} = (\pi+\varphi)/2 + \pi$.

\subsubsection{Scale-factor evolution and the algebraic determination of \texorpdfstring{$w_a$}{wa}}
\label{sec:eft:wa_derivation}

The constant-coupling viscosity of the preceding section fixes $w_0$
algebraically but predicts a \emph{time-independent} effective equation of
state.  DESI~DR2 data require a time-varying component; within SSEE this must
arise from a scale-factor-dependent extension of the viscosity coefficient.

\paragraph{Structural requirements on $\zeta(a)$.}
Any admissible SSEE viscosity extension $\zeta(a)$ must satisfy three
independent constraints:
\begin{enumerate}[label=(\roman*)]
  \item \textbf{De~Sitter endpoint:} $\zeta(a)\to 0$ as $a\to 1$, so that no
        viscous term shifts the present-day equation of state away from the
        k-essence value $w_0 = -T_r/M_v$.
  \item \textbf{Dimensional consistency:} $[\zeta] = [\rho_{\mathrm{DE}}/H]$,
        constructed from the available dynamical quantities
        $\{\rho_{\mathrm{DE}},H,a\}$.
  \item \textbf{Algebraic closure:} the dimensionless amplitude must be a
        ratio of SSEE structural constants derived exclusively from
        $\{\varphi,\pi\}$.
\end{enumerate}
The unique admissible form satisfying (i) and (ii) at leading order in $(1-a)$
— the natural expansion around the de~Sitter endpoint — is
%
\begin{equation}
  \zeta(a) = \Gamma\,\frac{(1-a)\,\rho_{\mathrm{DE}}(a)}{H(a)},
  \label{eq:zeta_general}
\end{equation}
%
where $\Gamma$ is a dimensionless amplitude to be fixed by constraint (iii).

\paragraph{Algebraic determination of $\Gamma$.}
The k-essence sector has already employed the ratio $T_r/M_v$ to fix
$|w_0|$.  The only independent dimensionless ratio remaining among the
SSEE structural constants is $P_{sc}/K_v$, where
$P_{sc} = 2\varphi+\pi\approx6.378$ (Dynamical Evolution Scalar) and
$K_v = 2(\varphi+\pi)\approx9.519$ (Structural Constraint).
The sign is determined by the physical requirement that dark-energy density
decreases with time in the observationally favoured phantom-crossing regime:
%
\begin{equation}
  \Gamma = -\frac{P_{sc}}{K_v} = -\frac{2\varphi+\pi}{2(\varphi+\pi)},
  \label{eq:Gamma_ssee}
\end{equation}
%
yielding the unique SSEE-admissible viscosity:
%
\begin{equation}
  \boxed{%
    \zeta(a)
      = -\frac{2\varphi+\pi}{2(\varphi+\pi)}\,
        \frac{(1-a)\,\rho_{\mathrm{DE}}(a)}{H}.}
  \label{eq:zeta_ssee}
\end{equation}

\paragraph{Algebraic consistency with CPL.}
The CPL parametrisation $w(a) = w_0 + w_a(1-a)$ is the standard
model-independent description of slowly-varying dark energy; it serves here as
the observational language, not as a physical assumption of SSEE.  Adopting CPL
and substituting $\zeta(a)$ from Eq.~(\ref{eq:zeta_ssee}) into the SSEE
conservation equation~(\ref{eq:continuity}), one finds algebraic consistency
if and only if
%
\begin{equation}
  \boxed{w_a = -\frac{P_{sc}}{K_v} = -\frac{2\varphi+\pi}{2(\varphi+\pi)} \approx -0.6699.}
  \label{eq:wa_ssee}
\end{equation}
%
This is not a derived prediction in the sense that CPL is deduced from the
action; rather, it establishes that \emph{the SSEE viscosity uniquely fixes
the CPL coefficient $w_a$ from the algebraic structure of $\{\varphi,\pi\}$}.
The quantity $w_a$ is therefore not a free parameter of the model but a
computable consequence of the two-axiom system, in the same sense that $w_0$
is fixed by the k-essence exponent $n$.  At $a=1$ the viscous correction
vanishes ($\zeta\to 0$), recovering the pure k-essence regime $w = w_0$.

The implied dark-energy density evolution reads:
%
\begin{equation}
  \rho_{\mathrm{DE}}(a)
    = \rho_{\mathrm{DE},0}\,
      a^{-3(1+w_0+w_a)}\,e^{-3w_a(1-a)},
  \label{eq:rho_CPL}
\end{equation}
%
consistent with the CPL conservation equation
%
\begin{equation}
  \dot\rho_{\mathrm{DE}}
    + 3H\,\rho_{\mathrm{DE}}\,(1+w_0)
    = -3H\,w_a\,(1-a)\,\rho_{\mathrm{DE}},
  \label{eq:cont_CPL}
\end{equation}
%
with the effective viscous pressure
%
\begin{equation}
  \Pi(a)
    = w_a\,(1-a)\,\rho_{\mathrm{DE}}(a).
  \label{eq:Pi_evolving}
\end{equation}
%
At $a \ll 1$ the factor $(1-a)\to 1$ and $\rho_{\mathrm{DE}}$ redshifts
according to Eq.~(\ref{eq:rho_CPL}), so $\zeta$ remains finite and
well-behaved throughout cosmic history.

\paragraph{Boundary conditions.}
The full viscosity interpolates between two algebraically-fixed regimes:
%
\begin{align}
  a \to 0 \;(\text{early}): &\quad
    \zeta \to -\frac{2\varphi+\pi}{2(\varphi+\pi)}\,
              \frac{\rho_{\mathrm{DE,0}}\,a^{-3(1+w_0+w_a)}\,e^{-3w_a}}{H(a)},
  \label{eq:zeta_early}\\[4pt]
  a \to 1 \;(\text{today}): &\quad
    \zeta \to 0.
  \label{eq:zeta_today}
\end{align}
%
Both limits are free of divergences and require no new parameters.

The viscous pressure is therefore
%
\begin{equation}
  \Pi(a)
    = \underbrace{-\frac{3\,\mathrm{KAL}_{0}\,H^{2}}{8\pi G}}_{\text{background}}
      \underbrace{-\;3H\,w_a\,(1-a)\,\rho_{\mathrm{DE}}(a)}_{\text{CPL evolution}},
  \label{eq:Pi_total}
\end{equation}
%
where the first term is the constant-coupling Eckart pressure that sustains
$\ddot a > 0$ and the second term generates the running $w(a)$ observed by
DESI~DR2 \citep{AbdulKarim2025}.

% ─────────────────────────────────────────────────────────────────────────────
\subsection{Modified Friedmann Equations}
\label{sec:eft:friedmann}

Applying the principle of least action ($\delta S = 0$) to the
FLRW metric $ds^{2} = -dt^{2} + a^{2}(t)\,d\mathbf{x}^{2}$ yields
three governing equations for the SSEE background.

\paragraph{I. Energy equation:}
\begin{equation}
  H^{2}
    = \frac{8\pi G}{3}
      \bigl(\rho_{m} + \rho_{r} + \rho_{\mathrm{SSEE}}\bigr).
  \label{eq:Friedmann_H}
\end{equation}

\paragraph{II. Acceleration equation:}
\begin{equation}
  \frac{\ddot{a}}{a}
    = -\frac{4\pi G}{3}
      \!\left[\rho_{m} + 2\rho_{r}
        + \rho_{\mathrm{SSEE}}\!\left(1 + 3w_{0}\right)
        - \frac{9\,\mathrm{KAL}_{0}\,H^{2}}{8\pi G}
      \right].
  \label{eq:Friedmann_a}
\end{equation}

\paragraph{III. Conservation equation with dissipative source:}
\begin{equation}
  \dot{\rho}_{\mathrm{SSEE}}
    + 3H\,\rho_{\mathrm{SSEE}}\!\left(1 + w_{0}\right)
    = \frac{9\,\mathrm{KAL}_{0}\,H^{3}}{8\pi G}.
  \label{eq:continuity}
\end{equation}

\noindent\textbf{Derivation of the coefficient 9.}\par\smallskip
\noindent From Eq.~(\ref{eq:zeta}), the viscous pressure is
$\Pi = -3\mathrm{KAL}_{0}H^{2}/8\pi G$.
In the Raychaudhuri equation the effective pressure of the SSEE sector enters as
$p_{\mathrm{eff}} = w_{0}\rho_{\mathrm{SSEE}} + \Pi$, so the acceleration term becomes
%
\begin{align}
  -\frac{4\pi G}{3}\bigl(\rho_{\mathrm{SSEE}} + 3p_{\mathrm{eff}}\bigr)
  &= -\frac{4\pi G}{3}
     \!\left[\rho_{\mathrm{SSEE}}(1+3w_{0}) + 3\Pi\right] \notag\\
  &= -\frac{4\pi G}{3}
     \!\left[\rho_{\mathrm{SSEE}}(1+3w_{0})
             - \frac{9\,\mathrm{KAL}_{0}\,H^{2}}{8\pi G}\right],
  \label{eq:deriv_coeff9}
\end{align}
which is exactly Eq.~(\ref{eq:Friedmann_a}).  The factor 9 arises from
$3 \times |\Pi|/(H^{2}) = 3 \times 3\mathrm{KAL}_{0}/8\pi G$; a factor
of 3 rather than 9 would correspond to including $\Pi$ only once instead
of as the trace of the viscous stress.

Similarly, the conservation law $\nabla_{\!\mu}T^{\mu\nu}=0$ for the effective
fluid $(w_{0},\Pi)$ gives
$\dot{\rho} + 3H(\rho + w_{0}\rho + \Pi) = 0$, hence
$\dot{\rho} + 3H\rho(1+w_{0}) = 9\mathrm{KAL}_{0}H^{3}/8\pi G$,
confirming that Eq.~(\ref{eq:Friedmann_a}) and Eq.~(\ref{eq:continuity})
are mutually consistent.

% ─────────────────────────────────────────────────────────────────────────────
\subsection{Israel--Stewart Viscous Perturbations}
\label{sec:eft:IS}

The Eckart formulation used in \S\ref{sec:eft:PX}--\ref{sec:eft:friedmann}
is a first-order theory that violates causality \citep{HiscockLindblom1985}:
the viscous signal propagates instantaneously.  Israel--Stewart (IS) theory
\citep{IsraelStewart1979} resolves this by promoting $\Pi$ to a dynamical
variable with a finite relaxation time $\tau_{\Pi}$:
%
\begin{equation}
  \tau_{\Pi} H_{0}\,\tilde{h}(a)\,
    \frac{\mathrm{d}\tilde{\Pi}}{\mathrm{d}\ln a}
  + \tilde{\Pi}
  = -\,\mathrm{KAL}_{0}\,\tilde{h}^{2}(a),
  \label{eq:IS_transport}
\end{equation}
%
where $\tilde{h} = H/H_{0}$, $\tilde{\Pi} = \Pi\,8\pi G/H_{0}^{2}$, and
all densities are normalised to $\rho_{\mathrm{crit},0}$.  In the
limit $\tau_{\Pi} \to 0$, Eq.~(\ref{eq:IS_transport}) reduces to the
Eckart result $\tilde{\Pi} = -\mathrm{KAL}_{0}\tilde{h}^{2}$.

\paragraph{Causality constraint.}
The bulk-viscous sound speed in IS theory is \citep{HiscockLindblom1985}
%
\begin{equation}
  c_{s}^{2}
    = \frac{\zeta}{\rho_{\mathrm{DE}}\,\tau_{\Pi}}
    = \frac{\mathrm{KAL}_{0}}{3\,\Omega_{\mathrm{DE}}\,(\tau_{\Pi} H_{0})}
    \leq 1.
  \label{eq:IS_causality}
\end{equation}
%
Setting $c_{s}^{2} = 1$ (the causality boundary) uniquely fixes the
IS relaxation time as an algebraic combination of SSEE constants:
%
\begin{equation}
  \tau_{\Pi} H_{0}
    \;=\; \frac{\mathrm{KAL}_{0}}{3\,\Omega_{\mathrm{DE}}}
    \;=\; \frac{\mathrm{KAL}_{0}\,M_{v}}{3\,T_{r}}
    \;\approx\; 2.191,
  \label{eq:tau_IS}
\end{equation}
%
with no free parameters.

\paragraph{Growth index.}
Coupling Eq.~(\ref{eq:IS_transport}) to the conservation
equation~(\ref{eq:continuity}) and the linear-growth ODE
%
\begin{equation}
  \delta'' + \left[2 + \frac{\mathrm{d}\ln\tilde{h}}{\mathrm{d}\ln a}\right]\delta'
  = \frac{3}{2}\,\frac{\Omega_{m,\mathrm{dyn}}\,a^{-3}}{\tilde{h}^{2}(a)}\,\delta,
  \label{eq:growth_IS}
\end{equation}
%
we solve the system numerically (\texttt{ssee\_is\_growth.py}) with
the IS background $\tilde{h}^{2}(a) = \Omega_{m,\mathrm{dyn}}\,a^{-3}+\rho_{\mathrm{DE}}(a)$.
The growth rate $f \equiv \mathrm{d}\ln\delta/\mathrm{d}\ln a \approx \Omega_{m}(a)^{\gamma}$
is well described by a power law over $0.1 \leq a \leq 1$ with
%
\begin{equation}
  \gamma_{\mathrm{IS}} = 0.657 \pm 0.002
  \quad(\text{essentially independent of }\tau_{\Pi} H_{0}).
  \label{eq:gamma_IS}
\end{equation}
%
This differs by 6\% from the algebraic sovereignty $\gamma = \varphi^{-1} = 0.618$.
Paper~3 \S5.4 reports both values: the algebraic prediction $\gamma_\mathrm{alg}=0.618$
and the IS numerical result $\gamma_\mathrm{IS}=0.657$.
The IS formalism provides the causal theoretical structure; whether $\gamma = \varphi^{-1}$
is the true attractor of the IS boundary conditions remains a target prediction
for forthcoming Stage-IV weak-lensing surveys (LSST, Euclid).
Note that the direct IS growth-factor integral (this Appendix) yields $S_8=0.837$,
while Paper~3's CAMB post-processing method with the same $\gamma_\mathrm{IS}=0.657$
gives $S_8=0.818$; both are consistent with Planck 2018 ($0.832\pm0.013$) within $1\sigma$.

\paragraph{$S_{8}$ constraint.}
Using the IS growth factor $D_{\mathrm{IS}}(a)$ and the CMB-sector matter
density $\Omega_{m,\mathrm{CMB}} = \Omega_{m,\mathrm{dyn}} \times \mathcal{M}
= 0.3199$ (MIRA, Paper~3),
%
\begin{equation}
  S_{8}^{\mathrm{IS}}
    \equiv \sigma_{8,\mathrm{IS}}
           \sqrt{\frac{\Omega_{m,\mathrm{CMB}}}{0.3}}
    \approx 0.837,
  \label{eq:S8_IS}
\end{equation}
%
within $0.4\sigma$ of the Planck-2018 value $0.832 \pm 0.013$~\citep{PlanckVI2020}.
Using the dynamic sector $\Omega_{m,\mathrm{dyn}} = 0.160$ instead gives
$S_{8} \approx 0.59$ ($6\sigma$ tension); the correct physical choice depends
on the formal resolution of the two-sector Friedmann problem (Section~\ref{sec:eft:status}, item 3).

% ─────────────────────────────────────────────────────────────────────────────
\subsection{Physical Regimes}
\label{sec:eft:regimes}

\paragraph{Late universe ($z \lesssim 2$, DESI~DR2).}
At low redshift $H \approx H_{0}$, so the viscous source
$\propto \mathrm{KAL}_{0}H^{3}/8\pi G$ is small compared to
$\rho_{\mathrm{SSEE}}$.  The acceleration equation is dominated by
$\rho_{\mathrm{SSEE}}(1+3w_{0}) < 0$
(since $w_{0} \approx -0.840 < -1/3$), driving $\ddot{a} > 0$
without dark matter.  This regime is the one tested against DESI~DR2
in Paper~2, where the algebraic prediction $(w_{0},w_{a})$ matches
the data at $0.05\sigma$.

\paragraph{Early universe (CMB epoch, $z \sim 10^{3}$).}
At recombination $H \gg H_{0}$, the geometric retention term
$9\,\mathrm{KAL}_{0}H^{3}/8\pi G$ acts as a dissipative source in
Eq.~(\ref{eq:continuity}), injecting energy into the dark-energy fluid
and modifying its effective equation of state.  This structural
viscosity provides the physical mechanism by which an effective
matter density $\Omega_{m,\mathrm{eff}} = 1 - T_{r}/M_{v} \approx 0.160$
yields the correct CMB acoustic scale through the MIRA factor
$\mathcal{M} = 2 - 1/\varphi^{2} \approx 1.9989$ (Paper~3, \S2.3).

% ─────────────────────────────────────────────────────────────────────────────
% ─────────────────────────────────────────────────────────────────────────────
\subsection{Inflationary Embedding: \texorpdfstring{$\alpha$}{α}-Attractor with \texorpdfstring{$\alpha=\varphi^{4}/3$}{α=φ⁴/3}}
\label{sec:eft:inflation}

The SSEE algebraic predictions for the primordial scalar spectral index
$n_s = 1 - \varphi^{-7}$ (Paper~1, \S\ref{sec:register}) and the
tensor-to-scalar ratio $r = \varphi^{-10}$ (Paper~4) belong simultaneously
to a well-defined class of inflationary models: the
\emph{$\alpha$-attractors} of Kallosh \& Linde \citeyearpar{KalloshLinde2013}.

\paragraph{$\alpha$-attractor predictions.}
In the universal leading-order limit, $\alpha$-attractor models predict
\begin{equation}
  n_s = 1 - \frac{2}{N}, \qquad
  r   = \frac{12\alpha}{N^{2}},
  \label{eq:alpha_attractor}
\end{equation}
where $N$ is the number of inflationary e-folds and $\alpha > 0$
parameterises the inverse curvature of the scalar field's K\"{a}hler target
manifold \citep{KalloshLinde2013,Ferrara2013}.
Setting $\alpha = 1$ recovers Starobinsky $R^{2}$ inflation exactly.

\paragraph{SSEE identification.}
Inserting $n_s = 1 - \varphi^{-7}$ into the first relation gives
\begin{equation}
  N = 2/\varphi^{-7} = 2\varphi^{7} \approx 58.07,
  \label{eq:N_ssee}
\end{equation}
which lies squarely in the observationally required window $50\lesssim N\lesssim 60$.
Substituting $N = 2\varphi^{7}$ and $r = \varphi^{-10}$ into the second relation:
\begin{equation}
  \alpha = \frac{r\,N^{2}}{12}
         = \frac{\varphi^{-10}\cdot 4\varphi^{14}}{12}
         = \frac{4\varphi^{4}}{12}
         = \frac{\varphi^{4}}{3}.
  \label{eq:alpha_ssee}
\end{equation}
This result is algebraically exact (verified numerically to machine precision;
see \texttt{ssee\_inflation\_connection.py}).
Using the Fibonacci identity $\varphi^{4} = 3\varphi + 2$:
\begin{equation}
  \alpha = \frac{3\varphi + 2}{3} = \varphi + \frac{2}{3} \approx 2.285.
  \label{eq:alpha_fibonacci}
\end{equation}

\paragraph{Geometric interpretation.}
In the $\alpha$-attractor framework the K\"{a}hler curvature of the
inflaton's target space (a Poincar\'{e} disk of radius $\sqrt{3\alpha}$)
is $\mathcal{R} = -2/(3\alpha)$.
For $\alpha = \varphi^{4}/3$:
\begin{equation}
  \mathcal{R} = -\frac{2}{3 \cdot \varphi^{4}/3} = -\frac{2}{\varphi^{4}}
               = -\varphi^{-4} \approx -0.146.
  \label{eq:kahler_curvature}
\end{equation}
The inflaton therefore lives on a hyperbolic manifold whose curvature is
set by $\varphi$ alone. This gives a \emph{geometric} origin for the
appearance of $\varphi$ in inflationary observables: $\varphi$ enters not
as a numerical coincidence but as the intrinsic length scale of the
K\"{a}hler geometry of the SSEE inflaton sector.

\paragraph{Slow-roll parameters.}
The SSEE predictions $(n_s,\,r) = (1-\varphi^{-7},\,\varphi^{-10})$ fix the
slow-roll parameters uniquely:
\begin{align}
  \varepsilon &= \frac{r}{16} = \frac{\varphi^{-10}}{16}, \label{eq:sr_eps} \\
  \eta &= \frac{n_s - 1 + 6\varepsilon}{2}
        = \frac{-\varphi^{-7} + 3\varphi^{-10}/8}{2}, \label{eq:sr_eta} \\
  \frac{\eta}{\varepsilon} &= 3 - 8\varphi^{3} = -5 - 16\varphi \approx -30.9.
  \label{eq:sr_ratio}
\end{align}
The ratio $\eta/\varepsilon = 3 - 8\varphi^{3}$ is a pure algebraic
identity derived from the SSEE axioms $\{\varphi, \pi\}$; it defines a
one-parameter family of inflationary potentials $V(\phi_{\rm inf})$ whose
single element consistent with $\alpha$-attractors is $\alpha = \varphi^{4}/3$.

\paragraph{Relation to Starobinsky.}
Starobinsky ($\alpha = 1$) reproduces $n_s$ correctly for $N = 2\varphi^{7}$
but gives $r_{\rm Staro} = 3/\varphi^{14} \approx 0.00356$, a factor of
$\varphi^{4}/3\approx 2.28$ below the SSEE prediction $\varphi^{-10}$.
SSEE is therefore \emph{not} Starobinsky inflation; it is a generalised
$\alpha$-attractor that reduces to Starobinsky only in the limit
$\varphi \to 1$ (i.e.\ $\alpha \to 1$, the fixed point of the golden-ratio
recursion $x = 1 + 1/x$).

% ─────────────────────────────────────────────────────────────────────────────
\subsection{Summary and Open Problems}
\label{sec:eft:status}

Table~\ref{tab:eft_summary} summarises the EFT status.
The action~(\ref{eq:eft_action}) with kinetic exponent~(\ref{eq:n_ssee})
and viscosity~(\ref{eq:zeta_ssee}) constitutes a \emph{complete},
zero-free-parameter EFT description of the SSEE background evolution.
Every constant ($n$, $\mathrm{KAL}_{0}$, $w_0$, $w_a$, $\Omega_{\mathrm{DE}}$) is
derived from $\varphi$ and $\pi$ alone; only the overall energy scale $H_{0}$
enters as an observational input.  Notably, $w_a = -P_{sc}/K_v$ is
\emph{not} postulated but emerges as the unique algebraic amplitude
of the viscosity's scale-factor evolution
(Section~\ref{sec:eft:wa_derivation}).

Four aspects require further development:
%
\begin{enumerate}
  \item \textbf{Growth index.}
        The IS-derived growth index $\gamma_{\mathrm{IS}} = 0.657$
        (\S\ref{sec:eft:IS}) differs by 6\% from the algebraic prediction
        $\varphi^{-1} = 0.618$.  A mechanism that would select
        $\gamma = \varphi^{-1}$ from IS boundary conditions is not yet
        identified; this constitutes a near-prediction testable by
        LSST/Euclid Stage-IV weak lensing.

  \item \textbf{Full perturbation spectrum.}
        The IS growth ODE~(\ref{eq:growth_IS}) governs sub-Hubble dark-matter
        perturbations; the matching to CMB scalar transfer functions and
        the tensor-mode propagation in the IS background require a complete
        Boltzmann implementation.

  \item \textbf{Radiative stability.}
        The EFT cutoff scale $\Lambda_{\mathrm{UV}}$ and the quantum stability
        of the $n < 0$ power-law sector are not yet established.

  \item \textbf{Inflation--dark energy unification.}
        The $\alpha$-attractor sector (\S\ref{sec:eft:inflation}) and the
        k-essence dark-energy sector (\S\ref{sec:eft:PX}) both originate
        from $\varphi$, suggesting a single unified action.
        Constructing the quintessential inflation potential $V(\phi)$
        that reduces to $P(X) = \mathcal{A}X^n$ at late times
        and to an $\alpha=\varphi^4/3$ attractor at early times
        is identified as high-priority future work.
\end{enumerate}

\begin{table}[h]
\centering
\caption{EFT status summary for SSEE-V3.6.}
\label{tab:eft_summary}
\begin{tabular}{lll}
\hline\hline
Element & Status & Derivation \\
\hline
Action form $S$ & \checkmark Complete & Standard GR + k-essence + Eckart \\
$u_{\mu}(\phi)$ coupling & \checkmark Complete & $u_{\mu} = -\partial_{\mu}\phi/\sqrt{2X}$ \\
Kinetic exponent $n$ & \checkmark Algebraic & $n=(T_{r}-M_{v})/2T_{r}$ from $\varphi,\pi$ \\
Equation of state $w_{0}$ & \checkmark Algebraic & $w_{0}=1/(2n-1)=-T_{r}/M_{v}$ \\
Viscosity coupling $\zeta$ & \checkmark Algebraic & $\zeta=\mathrm{KAL}_{0}H/8\pi G$ \\
Friedmann eqs.\ I--III & \checkmark Consistent & Factor 9 verified in \S\ref{sec:eft:friedmann} \\
$w_{a}$ from EFT & \checkmark Algebraic & $\zeta(a)\!\propto\!(1{-}a)\rho_\mathrm{DE}/H$; see \S\ref{sec:eft:wa_derivation} \\
IS relaxation time $\tau_\Pi$ & \checkmark Algebraic & $\tau_\Pi H_0 = \mathrm{KAL}_0/(3\Omega_\mathrm{DE}) \approx 2.191$ \\
IS causality ($c_s^2 \leq 1$) & \checkmark Verified & Eq.~(\ref{eq:IS_causality}); $c_s^2 = 1$ at boundary \\
Growth index $\gamma_\mathrm{IS}$ & \checkmark Numerical & $0.657 \pm 0.002$ (6\% above $\varphi^{-1}=0.618$) \\
$S_8$ (CMB sector) & \checkmark Numerical & $0.837$, within $0.4\sigma$ of Planck \\
$\alpha$-attractor ($\alpha=\varphi^4/3$) & \checkmark Algebraic & Eq.~(\ref{eq:alpha_ssee}); exact, $N=2\varphi^7$ \\
K\"{a}hler curvature $\mathcal{R}=-\varphi^{-4}$ & \checkmark Algebraic & Eq.~(\ref{eq:kahler_curvature}) \\
Full perturbation spectrum & $\circ$ Open & Off-diagonal CMB cov.\ + Stage-IV WL \\
Inflation--DE unification & $\circ$ Open & Quintessential inflation potential \\
Radiative stability & $\circ$ Open & Future work \\
\hline
\end{tabular}
\end{table}
